\section{Các phương pháp thể tích hữu hạn}
\subsection{Phương pháp thể tích hữu hạn của định luật bảo toàn}
Trước tiên, ta có dạng tổng quát của phương trình Euler cho trường hợp một chiều như sau:
\[
w_t + f(w)_x = 0.
\]
Ta có thể diễn đạt hệ phương trình này thành các phương trình khác nhau tương ứng với:
\begin{itemize}
    \item \textbf{Bảo toàn khối lượng:}
    \[
    \dfrac{\partial \rho}{\partial t} + \dfrac{\partial (\rho u)}{\partial x} = 0.
    \]
    \item \textbf{Bảo toàn động lượng:}
    \[
    \dfrac{\partial (\rho u)}{\partial t} + \dfrac{\partial (\rho u^2 + p)}{\partial x} = 0.
    \]
    \item \textbf{Bảo toàn năng lượng:}
    \[
    \dfrac{\partial E}{\partial t} + \dfrac{\partial (u(E + p))}{\partial x} = 0.
    \]
\end{itemize}

Ta xét hệ phương trình trên trên một miền tính toán là $\Omega = \{(x, t) : 0 \leq t \leq t_{\text{final}}, a(t) \leq x \leq b(t) \}$, với $t_{\text{final}}$ là thời điểm cuối cùng mà ta muốn khảo sát, $a(t)$, $b(t)$ là biên của miền tính toán, thỏa điều kiện đầu $u(x, 0) = u_0(x)$.

Ta chia nhỏ đoạn $[a, b]$ thành $N$ đoạn nhỏ hơn, ký hiệu đoạn con thứ $i$ là $[x_{i - 1/2}, x_{i + 1/2}]$, $i = \overline{1, ..., N}$. Tiếp theo, ta sẽ rời rạc hóa thời gian như sau: với giá trị $\Delta t$ cố định, ta chia nhỏ đoạn $[0, t_{\text{final}}]$ thành các đoạn nhỏ hơn với khoảng cách là $\Delta t$. Khi đó, miền $\Omega$ ban đầu sẽ được phân hoạch thành các miền nhỏ hơn: $[x_{i - 1/2}, x_{i + 1/2}] \times [t^n, t^{n+1}]$ ($t^n$ là mốc thời gian thứ $n$ khi chia nhỏ miền thời gian - xem ảnh minh họa bên dưới).


\begin{figure}[H]
\centering
\begin{tikzpicture}[scale=1.2]
    % Draw axes
    \draw[->] (0,0) -- (6,0) node[right] {$x$};
    \draw[->] (0,0) -- (0,4) node[above] {$t$};
    
    % Draw the domain cell with solid lines
    \draw[thick] (2,1.5) rectangle (4,2.5);
    
    % Dashed lines from axes to domain cell
    \draw[dashed] (2,0) -- (2,1.5);
    \draw[dashed] (4,0) -- (4,1.5);
    \draw[dashed] (0,1.5) -- (2,1.5);
    \draw[dashed] (0,2.5) -- (2,2.5);
    
    % Tick marks on x-axis
    \draw (2,0) -- (2,-0.1);
    \draw (4,0) -- (4,-0.1);
    
    % Tick marks on t-axis
    \draw (0,1.5) -- (-0.1,1.5);
    \draw (0,2.5) -- (-0.1,2.5);
    
    % x-axis labels
    \node[below] at (2,0) {$x_{i-1/2}$};
    \node[below] at (4,0) {$x_{i+1/2}$};
    
    % t-axis labels
    \node[left] at (0,1.5) {$t^n$};
    \node[left] at (0,2.5) {$t^{n+1}$};
    
    % Arrow pointing to the domain with label outside
    \draw[->] (5,3) -- (3.5,2.2);
    \node[right] at (5,3) {$[x_{i-1/2}, x_{i+1/2}] \times [t^n, t^{n+1}]$};
    
\end{tikzpicture}
\caption{Rời rạc hóa miền tính toán $\Omega$ thành các miền con}
\label{fig:domain_discretization}
\end{figure}

Lấy tích phân của phương trình Euler trên miền $[x_{i - 1/2}, x_{i + 1/2}] \times [t^n, t^{n+1}]$, ta có:
\begin{equation*}
    \begin{aligned}
        &\int_{t_n}^{t_{n+1}} \int_{x_{i-1/2}}^{x_{i+1/2}} \left( \frac{\partial w}{\partial t} + \frac{\partial f(w)}{\partial x} \right) dx \, dt = 0 \\
        \Leftrightarrow & \int_{t_n}^{t_{n+1}} \left( \int_{x_{i-1/2}}^{x_{i+1/2}}\frac{\partial w}{\partial t} dx \right) dt + \int_{t_n}^{t_{n+1}} \left( \int_{x_{i-1/2}}^{x_{i+1/2}} \frac{\partial f(w)}{\partial x} dx \right) dt = 0 \\
        \Leftrightarrow &\int_{x_{i-1/2}}^{x_{i+1/2}} \left( \int_{t_n}^{t_{n+1}} \frac{\partial w}{\partial t} dt \right) dx + \int_{t_n}^{t_{n+1}} \left( \int_{x_{i-1/2}}^{x_{i+1/2}} \frac{\partial f(w)}{\partial x} dx \right) dt = 0.
    \end{aligned}
\end{equation*}
Áp dụng Định lý cơ bản của giải tích (Newton-Leibniz):
\begin{itemize}
    \item Với biến thời gian $t$:
    $$ \int_{t_n}^{t_{n+1}} \frac{\partial w(x,t)}{\partial t} dt = w(x, t_{n+1}) - w(x, t_n) $$
    \item Với biến không gian $x$:
    $$ \int_{x_{i-1/2}}^{x_{i+1/2}} \frac{\partial f(w(x,t))}{\partial x} dx = f(w(x_{i+1/2}, t)) - f(w(x_{i-1/2}, t)) $$
\end{itemize}
Thay kết quả trên vào phương trình tích phân ban đầu, ta thu được phương trình cân bằng trên miền con $[x_{i - 1/2}, x_{i + 1/2}] \times [t^n, t^{n+1}]$:
\begin{equation*}
    \begin{aligned}
        &\int_{x_{i-1/2}}^{x_{i+1/2}} \left( \int_{t_n}^{t_{n+1}} \frac{\partial w}{\partial t} dt \right) dx + \int_{t_n}^{t_{n+1}} \left( \int_{x_{i-1/2}}^{x_{i+1/2}} \frac{\partial f(w)}{\partial x} dx \right) dt = 0 \\
       \Leftrightarrow &\int_{x_{i-1/2}}^{x_{i+1/2}} \left( w(x, t_{n+1}) - w(x, t_n)\right) dx + \int_{t_n}^{t_{n+1}} \left( f(w(x_{i+1/2}, t)) - f(w(x_{i-1/2}, t))\right) dt = 0 \\
        \Leftrightarrow &\int_{x_{i-1/2}}^{x_{i+1/2}} w(x, t_{n+1}) dx - \int_{x_{i-1/2}}^{x_{i+1/2}} w(x, t_n) dx + \int_{t_n}^{t_{n+1}} f(w(x_{i+1/2}, t)) dt - \int_{t_n}^{t_{n+1}} f(w(x_{i-1/2}, t)) dt = 0
    \end{aligned}
\end{equation*}

Ta áp dụng các ký hiệu:
\[
[W]_i^n = \dfrac{1}{\Delta x_i}\int_{x_{i-1/2}}^{x_{i+1/2}} w(x, t_n) dx, \quad F_{i+1/2} = \dfrac{1}{\Delta t}\int_{t_n}^{t_{n+1}} f(w(x_{i+1/2}, t)) dt
\]
với $\Delta x_i = x_{i+1/2} - x_{i-1/2}$ và $\Delta t = t^{n+1} - t^n$. Bằng cách ký hiệu tương tự cho $[W]_i^{n+1}$ và $F_{i-1/2}$, ta có:
$$
\Delta x_i [W]_i^{n+1} - \Delta x_i [W]_i^n + \Delta t F_{i+1/2} - \Delta t F_{i-1/2} = 0.
$$

Chia cả hai vế phương trình cho $\Delta x_i$, ta thu được:
\begin{equation*}
    [W]_i^{n+1} = [W]_i^n - \frac{\Delta t}{\Delta x_i} \left( F_{i+1/2} - F_{i-1/2} \right).
\end{equation*}

Giả sử giá trị của $F^n_{i+1/2}$ chỉ phụ thuộc vào các giá trị của $W_i^n$ và $W_{i+1}^{n}$, khi đó ta có thể xấp xỉ $F^n_{i+1/2}$ như sau:

\[
F^n_{i+1/2} \approx \mathcal{F} (W_i^n, W^n_{i+1}).
\]
Bằng cách xấp xỉ tương tự cho $F^n_{i-1/2}$, ta có kết quả sau:
\[
[W]_i^{n+1} = [W]_i^n - \frac{\Delta t}{\Delta x_i} \left( \mathcal{F} (W_i^n, W^n_{i+1}) - \mathcal{F} (W_{i-1}^n, W^n_{i})\right).
\]

Cách xấp xỉ như trên làm cho $W_i^{n+1}$ chỉ phụ thuộc vào $W_{i-1}^{n}, W_i^{n}$ và $W_{i+1}^{n}$, do đó ta gọi nó là phương pháp three-point stencil. Với mỗi cách chọn $\mathcal{F}$, ta có được một phương pháp số cho phương trình Euler.

Ta chọn $F^n_{i+1/2} \approx \mathcal{F} (W_i^n, W^n_{i+1}) = \dfrac{1}{2}(f(W_i^n) + f(W_{i+1}^n))$, khi đó phương trình cập nhật trở thành:
\begin{align*}
    W_i^{n+1} &= W_i^n - \frac{\Delta t}{\Delta x} \left[ \underbrace{\frac{1}{2}(f(W_i^n) + f(W_{i+1}^n))}_{F_{i+1/2}} - \underbrace{\frac{1}{2}(f(W_{i-1}^n) + f(W_i^n))}_{F_{i-1/2}} \right] \\
    &= W_i^n - \frac{\Delta t}{2\Delta x} \left[ f(W_i^n) + f(W_{i+1}^n) - f(W_{i-1}^n) - f(W_i^n) \right]
\end{align*}
Các số hạng chứa $f(W_i^n)$ triệt tiêu lẫn nhau, ta thu được phương trình sai phân trung tâm (Central Difference Scheme):
\begin{equation*}
    W_i^{n+1} = W_i^n - \frac{\Delta t}{2\Delta x} \left( f(W_{i+1}^n) - f(W_{i-1}^n) \right).
\end{equation*}
\textbf{Nhận xét:} Mặc dù phương pháp này đơn giản và trực quan, nhưng nó được chứng minh là \textbf{không ổn định (unstable)} đối với các bài toán hyperbolic (như hệ phương trình Euler) nếu không kèm theo các điều kiện rất nghiêm ngặt. Do đó, ta cần sử dụng phương pháp thay thế ổn định hơn, đó là phương pháp Lax-Friedrichs.
\subsection{Phương pháp Lax-Friedrich và local Lax-Friedrich}
\subsubsection{Lax-Friedrich}
Phương pháp Lax-Friedrich là một trong những phương pháp đơn giản nhất để giải hệ phương trình Euler. Bằng cách thay đổi $W_i^n = \dfrac{W_{i+1}^n + W_{i-1}^n}{2}$ vào phương trình sai phân trung tâm, ta có được phương pháp Lax-Friedrich:
\[
W_i^{n+1} = \dfrac{W_{i-1}^n + W_{i+1}^n}{2} - \frac{\Delta t}{2\Delta x} \left( f(W_{i+1}^n) - f(W_{i-1}^n) \right).
\]

Dưới đây là kết quả áp dụng phương pháp Lax-Friedrich cho bài toán ống sốc Sod:

\begin{figure}[H]
\centering
\begin{tikzpicture}
    % Define spacing and dimensions
    \def\plotwidth{7cm}
    \def\plotheight{5.5cm}
    \def\hspace{8cm}
    \def\vspace{6.5cm}
    
    % Top row
    \node[anchor=south west, fill=white] at (0,\vspace) {
        \resizebox{\plotwidth}{\plotheight}{% This file was created by matlab2tikz.
%
%The latest updates can be retrieved from
%  http://www.mathworks.com/matlabcentral/fileexchange/22022-matlab2tikz-matlab2tikz
%where you can also make suggestions and rate matlab2tikz.
%
\definecolor{mycolor1}{rgb}{0.12941,0.12941,0.12941}%
%
\begin{tikzpicture}

\begin{axis}[%
width=3.229in,
height=2.498in,
at={(0.542in,0.392in)},
scale only axis,
xmin=0,
xmax=1,
xlabel style={font=\color{mycolor1}},
xlabel={x},
ymin=0.1,
ymax=1,
ylabel style={font=\color{mycolor1}},
ylabel={$\text{Density (}\rho\text{)}$},
axis background/.style={fill=white},
title style={font=\bfseries\color{mycolor1}},
title={Lax-Friedrichs Density Distribution},
xmajorgrids,
ymajorgrids,
legend style={legend cell align=left, align=left}
]
\addplot [color=blue, line width=2.0pt]
  table[row sep=crcr]{%
0	0.999998258674086\\
0.0050251256281407	0.999998258674086\\
0.0100502512562814	0.999997158031474\\
0.0150753768844221	0.999997158031474\\
0.0201005025125628	0.999993960385556\\
0.0251256281407035	0.999993960385556\\
0.0301507537688442	0.999987331053824\\
0.0351758793969849	0.999987331053824\\
0.0402010050251256	0.999974306385863\\
0.0452261306532663	0.999974306385863\\
0.050251256281407	0.999949649995506\\
0.0552763819095477	0.999949649995506\\
0.0603015075376884	0.999904600941577\\
0.0653266331658292	0.999904600941577\\
0.0703517587939698	0.999825110253605\\
0.0753768844221105	0.999825110253605\\
0.0804020100502513	0.999689571424351\\
0.085427135678392	0.999689571424351\\
0.0904522613065327	0.999466129261941\\
0.0954773869346734	0.999466129261941\\
0.100502512562814	0.999109792029165\\
0.105527638190955	0.999109792029165\\
0.110552763819095	0.998559738035543\\
0.115577889447236	0.998559738035543\\
0.120603015075377	0.997737356666548\\
0.125628140703518	0.997737356666548\\
0.130653266331658	0.996545638678143\\
0.135678391959799	0.996545638678143\\
0.14070351758794	0.994870475887804\\
0.14572864321608	0.994870475887804\\
0.150753768844221	0.992584215482233\\
0.155778894472362	0.992584215482233\\
0.160804020100503	0.989551455854656\\
0.165829145728643	0.989551455854656\\
0.170854271356784	0.985636641408037\\
0.175879396984925	0.985636641408037\\
0.180904522613065	0.980712624638726\\
0.185929648241206	0.980712624638726\\
0.190954773869347	0.974669128140635\\
0.195979899497487	0.974669128140635\\
0.201005025125628	0.967420028086455\\
0.206030150753769	0.967420028086455\\
0.21105527638191	0.958908595520745\\
0.21608040201005	0.958908595520745\\
0.221105527638191	0.949110205482028\\
0.226130653266332	0.949110205482028\\
0.231155778894472	0.938032450461357\\
0.236180904522613	0.938032450461357\\
0.241206030150754	0.925712969257643\\
0.246231155778894	0.925712969257643\\
0.251256281407035	0.912215554136316\\
0.256281407035176	0.912215554136316\\
0.261306532663317	0.897625204731094\\
0.266331658291457	0.897625204731094\\
0.271356783919598	0.882042773741923\\
0.276381909547739	0.882042773741923\\
0.281407035175879	0.865579738320631\\
0.28643216080402	0.865579738320631\\
0.291457286432161	0.848353478530375\\
0.296482412060302	0.848353478530375\\
0.301507537688442	0.830483289023292\\
0.306532663316583	0.830483289023292\\
0.311557788944724	0.812087217074702\\
0.316582914572864	0.812087217074702\\
0.321608040201005	0.793279720757579\\
0.326633165829146	0.793279720757579\\
0.331658291457286	0.774170076505514\\
0.336683417085427	0.774170076505514\\
0.341708542713568	0.754861430669598\\
0.346733668341709	0.754861430669598\\
0.351758793969849	0.735450377453233\\
0.35678391959799	0.735450377453233\\
0.361809045226131	0.716026948183909\\
0.366834170854271	0.716026948183909\\
0.371859296482412	0.696674907741938\\
0.376884422110553	0.696674907741938\\
0.381909547738693	0.67747226815296\\
0.386934673366834	0.67747226815296\\
0.391959798994975	0.658491943391131\\
0.396984924623116	0.658491943391131\\
0.402010050251256	0.639802481053911\\
0.407035175879397	0.639802481053911\\
0.412060301507538	0.621468814323378\\
0.417085427135678	0.621468814323378\\
0.422110552763819	0.603552980638621\\
0.42713567839196	0.603552980638621\\
0.4321608040201	0.586114751247608\\
0.437185929648241	0.586114751247608\\
0.442211055276382	0.56921210808866\\
0.447236180904523	0.56921210808866\\
0.452261306532663	0.552901491547567\\
0.457286432160804	0.552901491547567\\
0.462311557788945	0.537237725666756\\
0.467336683417085	0.537237725666756\\
0.472361809045226	0.522273508934371\\
0.477386934673367	0.522273508934371\\
0.482412060301508	0.508058343748125\\
0.487437185929648	0.508058343748125\\
0.492462311557789	0.494636774126054\\
0.49748743718593	0.494636774126054\\
0.50251256281407	0.482045821043267\\
0.507537688442211	0.482045821043267\\
0.512562814070352	0.470311562847773\\
0.517587939698492	0.470311562847773\\
0.522613065326633	0.459444919804426\\
0.527638190954774	0.459444919804426\\
0.532663316582915	0.4494368765261\\
0.537688442211055	0.4494368765261\\
0.542713567839196	0.4402536086084\\
0.547738693467337	0.4402536086084\\
0.552763819095477	0.431832239809118\\
0.557788944723618	0.431832239809118\\
0.562814070351759	0.424078182011587\\
0.5678391959799	0.424078182011587\\
0.57286432160804	0.416865114518571\\
0.577889447236181	0.416865114518571\\
0.582914572864322	0.410038550907304\\
0.587939698492462	0.410038550907304\\
0.592964824120603	0.40342356319841\\
0.597989949748744	0.40342356319841\\
0.603015075376884	0.396836598699914\\
0.608040201005025	0.396836598699914\\
0.613065326633166	0.390100539058563\\
0.618090452261307	0.390100539058563\\
0.623115577889447	0.383061390644362\\
0.628140703517588	0.383061390644362\\
0.633165829145729	0.375604456829846\\
0.638190954773869	0.375604456829846\\
0.64321608040201	0.367667676009173\\
0.648241206030151	0.367667676009173\\
0.653266331658292	0.359250069152967\\
0.658291457286432	0.359250069152967\\
0.663316582914573	0.350413875709734\\
0.668341708542714	0.350413875709734\\
0.673366834170854	0.341279834325959\\
0.678391959798995	0.341279834325959\\
0.683417085427136	0.332016016323179\\
0.688442211055276	0.332016016323179\\
0.693467336683417	0.322821481045337\\
0.698492462311558	0.322821481045337\\
0.703517587939699	0.313906658272723\\
0.708542713567839	0.313906658272723\\
0.71356783919598	0.305472675771766\\
0.718592964824121	0.305472675771766\\
0.723618090452261	0.297691776496801\\
0.728643216080402	0.297691776496801\\
0.733668341708543	0.290690479155551\\
0.738693467336683	0.290690479155551\\
0.743718592964824	0.284536229357029\\
0.748743718592965	0.284536229357029\\
0.753768844221106	0.279227002181657\\
0.758793969849246	0.279227002181657\\
0.763819095477387	0.274681730293125\\
0.768844221105528	0.274681730293125\\
0.773869346733668	0.270727718442816\\
0.778894472361809	0.270727718442816\\
0.78391959798995	0.267079824438749\\
0.78894472361809	0.267079824438749\\
0.793969849246231	0.263306442201109\\
0.798994974874372	0.263306442201109\\
0.804020100502513	0.258782667092942\\
0.809045226130653	0.258782667092942\\
0.814070351758794	0.252649132000028\\
0.819095477386935	0.252649132000028\\
0.824120603015075	0.243837521263082\\
0.829145728643216	0.243837521263082\\
0.834170854271357	0.23128983941941\\
0.839195979899497	0.23128983941941\\
0.844221105527638	0.214515904876155\\
0.849246231155779	0.214515904876155\\
0.85427135678392	0.19439651193949\\
0.85929648241206	0.19439651193949\\
0.864321608040201	0.173548335473074\\
0.869346733668342	0.173548335473074\\
0.874371859296482	0.155354259314585\\
0.879396984924623	0.155354259314585\\
0.884422110552764	0.142055548427925\\
0.889447236180904	0.142055548427925\\
0.894472361809045	0.133771801376835\\
0.899497487437186	0.133771801376835\\
0.904522613065327	0.129224902756225\\
0.909547738693467	0.129224902756225\\
0.914572864321608	0.126943118672822\\
0.919597989949749	0.126943118672822\\
0.924623115577889	0.125864436203158\\
0.92964824120603	0.125864436203158\\
0.934673366834171	0.125374576407908\\
0.939698492462312	0.125374576407908\\
0.944723618090452	0.125158561602955\\
0.949748743718593	0.125158561602955\\
0.954773869346734	0.125065609579864\\
0.959798994974874	0.125065609579864\\
0.964824120603015	0.125026521322311\\
0.969849246231156	0.125026521322311\\
0.974874371859296	0.125010471494743\\
0.979899497487437	0.125010471494743\\
0.984924623115578	0.125004117809962\\
0.989949748743719	0.125004117809962\\
0.994974874371859	0.125002237860314\\
1	0.125002237860314\\
};
\addlegendentry{LF Solution}

\end{axis}
\end{tikzpicture}%}
    };
    
    \node[anchor=south west, fill=white] at (\hspace,\vspace) {
        \resizebox{\plotwidth}{\plotheight}{% This file was created by matlab2tikz.
%
%The latest updates can be retrieved from
%  http://www.mathworks.com/matlabcentral/fileexchange/22022-matlab2tikz-matlab2tikz
%where you can also make suggestions and rate matlab2tikz.
%
\definecolor{mycolor1}{rgb}{0.12941,0.12941,0.12941}%
%
\begin{tikzpicture}

\begin{axis}[%
width=3.229in,
height=2.498in,
at={(0.542in,0.392in)},
scale only axis,
xmin=0,
xmax=1,
xlabel style={font=\color{mycolor1}},
xlabel={x},
ymin=0,
ymax=1,
ylabel style={font=\color{mycolor1}},
ylabel={Velocity (u)},
axis background/.style={fill=white},
title style={font=\bfseries\color{mycolor1}},
title={Lax-Friedrichs Velocity Distribution},
xmajorgrids,
ymajorgrids,
legend style={legend cell align=left, align=left}
]
\addplot [color=red, line width=2.0pt]
  table[row sep=crcr]{%
0	2.06036324041918e-06\\
0.0050251256281407	2.06036324041918e-06\\
0.0100502512562814	3.36265861710774e-06\\
0.0150753768844221	3.36265861710774e-06\\
0.0201005025125628	7.14615245128612e-06\\
0.0251256281407035	7.14615245128612e-06\\
0.0301507537688442	1.49900430316176e-05\\
0.0351758793969849	1.49900430316176e-05\\
0.0402010050251256	3.04009126164035e-05\\
0.0452261306532663	3.04009126164035e-05\\
0.050251256281407	5.95744020942169e-05\\
0.0552763819095477	5.95744020942169e-05\\
0.0603015075376884	0.000112876346982274\\
0.0653266331658292	0.000112876346982274\\
0.0703517587939698	0.000206929492666156\\
0.0753768844221105	0.000206929492666156\\
0.0804020100502513	0.00036729980550932\\
0.085427135678392	0.00036729980550932\\
0.0904522613065327	0.000631684506043909\\
0.0954773869346734	0.000631684506043909\\
0.100502512562814	0.0010533406115355\\
0.105527638190955	0.0010533406115355\\
0.110552763819095	0.00170430169408887\\
0.115577889447236	0.00170430169408887\\
0.120603015075377	0.00267776348169563\\
0.125628140703518	0.00267776348169563\\
0.130653266331658	0.0040889420678909\\
0.135678391959799	0.0040889420678909\\
0.14070351758794	0.00607378414520779\\
0.14572864321608	0.00607378414520779\\
0.150753768844221	0.00878516643709386\\
0.155778894472362	0.00878516643709386\\
0.160804020100503	0.012386633240322\\
0.165829145728643	0.012386633240322\\
0.170854271356784	0.0170441958187177\\
0.175879396984925	0.0170441958187177\\
0.180904522613065	0.0229171259793809\\
0.185929648241206	0.0229171259793809\\
0.190954773869347	0.0301488978468447\\
0.195979899497487	0.0301488978468447\\
0.201005025125628	0.0388594026909874\\
0.206030150753769	0.0388594026909874\\
0.21105527638191	0.0491392987149729\\
0.21608040201005	0.0491392987149729\\
0.221105527638191	0.061046947689434\\
0.226130653266332	0.061046947689434\\
0.231155778894472	0.0746079510032226\\
0.236180904522613	0.0746079510032226\\
0.241206030150754	0.0898169335359206\\
0.246231155778894	0.0898169335359206\\
0.251256281407035	0.106640996031596\\
0.256281407035176	0.106640996031596\\
0.261306532663317	0.125024176695531\\
0.266331658291457	0.125024176695531\\
0.271356783919598	0.144892303814866\\
0.276381909547739	0.144892303814866\\
0.281407035175879	0.16615773836621\\
0.28643216080402	0.16615773836621\\
0.291457286432161	0.188723653543851\\
0.296482412060302	0.188723653543851\\
0.301507537688442	0.212487642212414\\
0.306532663316583	0.212487642212414\\
0.311557788944724	0.23734456283668\\
0.316582914572864	0.23734456283668\\
0.321608040201005	0.263188621392526\\
0.326633165829146	0.263188621392526\\
0.331658291457286	0.289914741780604\\
0.336683417085427	0.289914741780604\\
0.341708542713568	0.317419305777712\\
0.346733668341709	0.317419305777712\\
0.351758793969849	0.345600352689511\\
0.35678391959799	0.345600352689511\\
0.361809045226131	0.374357325522354\\
0.366834170854271	0.374357325522354\\
0.371859296482412	0.403590440478316\\
0.376884422110553	0.403590440478316\\
0.381909547738693	0.433199744397591\\
0.386934673366834	0.433199744397591\\
0.391959798994975	0.463083913811193\\
0.396984924623116	0.463083913811193\\
0.402010050251256	0.493138842149891\\
0.407035175879397	0.493138842149891\\
0.412060301507538	0.523256060631971\\
0.417085427135678	0.523256060631971\\
0.422110552763819	0.553321045611355\\
0.42713567839196	0.553321045611355\\
0.4321608040201	0.583211483006845\\
0.437185929648241	0.583211483006845\\
0.442211055276382	0.612795591187413\\
0.447236180904523	0.612795591187413\\
0.452261306532663	0.641930649267859\\
0.457286432160804	0.641930649267859\\
0.462311557788945	0.670461938612832\\
0.467336683417085	0.670461938612832\\
0.472361809045226	0.698222378575923\\
0.477386934673367	0.698222378575923\\
0.482412060301508	0.725033214173659\\
0.487437185929648	0.725033214173659\\
0.492462311557789	0.750706175171193\\
0.49748743718593	0.750706175171193\\
0.50251256281407	0.775047542408226\\
0.507537688442211	0.775047542408226\\
0.512562814070352	0.797864485629091\\
0.517587939698492	0.797864485629091\\
0.522613065326633	0.818973830208117\\
0.527638190954774	0.818973830208117\\
0.532663316582915	0.838213033475923\\
0.537688442211055	0.838213033475923\\
0.542713567839196	0.855452612225189\\
0.547738693467337	0.855452612225189\\
0.552763819095477	0.870608642655314\\
0.557788944723618	0.870608642655314\\
0.562814070351759	0.883653422502688\\
0.5678391959799	0.883653422502688\\
0.57286432160804	0.894622172394238\\
0.577889447236181	0.894622172394238\\
0.582914572864322	0.903613963632275\\
0.587939698492462	0.903613963632275\\
0.592964824120603	0.910785951841079\\
0.597989949748744	0.910785951841079\\
0.603015075376884	0.916341299503918\\
0.608040201005025	0.916341299503918\\
0.613065326633166	0.920512507653032\\
0.618090452261307	0.920512507653032\\
0.623115577889447	0.923542810537015\\
0.628140703517588	0.923542810537015\\
0.633165829145729	0.925668519424172\\
0.638190954773869	0.925668519424172\\
0.64321608040201	0.927104706127573\\
0.648241206030151	0.927104706127573\\
0.653266331658292	0.928035624273504\\
0.658291457286432	0.928035624273504\\
0.663316582914573	0.928610130715194\\
0.668341708542714	0.928610130715194\\
0.673366834170854	0.928941405250657\\
0.678391959798995	0.928941405250657\\
0.683417085427136	0.929109644715483\\
0.688442211055276	0.929109644715483\\
0.693467336683417	0.929166135978696\\
0.698492462311558	0.929166135978696\\
0.703517587939699	0.929137076640309\\
0.708542713567839	0.929137076640309\\
0.71356783919598	0.92902552716445\\
0.718592964824121	0.92902552716445\\
0.723618090452261	0.928809723491327\\
0.728643216080402	0.928809723491327\\
0.733668341708543	0.928435408298736\\
0.738693467336683	0.928435408298736\\
0.743718592964824	0.927798568433855\\
0.748743718592965	0.927798568433855\\
0.753768844221106	0.926712662535625\\
0.758793969849246	0.926712662535625\\
0.763819095477387	0.924850730793864\\
0.768844221105528	0.924850730793864\\
0.773869346733668	0.92164749765667\\
0.778894472361809	0.92164749765667\\
0.78391959798995	0.916140304151359\\
0.78894472361809	0.916140304151359\\
0.793969849246231	0.90672374526868\\
0.798994974874372	0.90672374526868\\
0.804020100502513	0.890802226908472\\
0.809045226130653	0.890802226908472\\
0.814070351758794	0.86437667993431\\
0.819095477386935	0.86437667993431\\
0.824120603015075	0.821759332524212\\
0.829145728643216	0.821759332524212\\
0.834170854271357	0.755961278742321\\
0.839195979899497	0.755961278742321\\
0.844221105527638	0.660804026508322\\
0.849246231155779	0.660804026508322\\
0.85427135678392	0.535792503536017\\
0.85929648241206	0.535792503536017\\
0.864321608040201	0.392343481928518\\
0.869346733668342	0.392343481928518\\
0.874371859296482	0.254309597865968\\
0.879396984924623	0.254309597865968\\
0.884422110552764	0.14571647010679\\
0.889447236180904	0.14571647010679\\
0.894472361809045	0.0753236230090875\\
0.899497487437186	0.0753236230090875\\
0.904522613065327	0.0361935796887442\\
0.909547738693467	0.0361935796887442\\
0.914572864321608	0.0165788812706796\\
0.919597989949749	0.0165788812706796\\
0.924623115577889	0.0073513466857154\\
0.92964824120603	0.0073513466857154\\
0.934673366834171	0.00317884511561765\\
0.939698492462312	0.00317884511561765\\
0.944723618090452	0.00134405423122103\\
0.949748743718593	0.00134405423122103\\
0.954773869346734	0.000555800041100139\\
0.959798994974874	0.000555800041100139\\
0.964824120603015	0.000224601651482036\\
0.969849246231156	0.000224601651482036\\
0.974874371859296	8.86670037435969e-05\\
0.979899497487437	8.86670037435969e-05\\
0.984924623115578	3.48649556972697e-05\\
0.989949748743719	3.48649556972697e-05\\
0.994974874371859	1.89471868760348e-05\\
1	1.89471868760348e-05\\
};
\addlegendentry{LF Solution}

\end{axis}
\end{tikzpicture}%}
    };
    
    % Bottom row
    \node[anchor=south west, fill=white] at (0,0) {
        \resizebox{\plotwidth}{\plotheight}{% This file was created by matlab2tikz.
%
%The latest updates can be retrieved from
%  http://www.mathworks.com/matlabcentral/fileexchange/22022-matlab2tikz-matlab2tikz
%where you can also make suggestions and rate matlab2tikz.
%
\definecolor{mycolor1}{rgb}{0.12941,0.12941,0.12941}%
%
\begin{tikzpicture}

\begin{axis}[%
width=3.229in,
height=2.498in,
at={(0.542in,0.392in)},
scale only axis,
xmin=0,
xmax=1,
xlabel style={font=\color{mycolor1}},
xlabel={x},
ymin=0.1,
ymax=1,
ylabel style={font=\color{mycolor1}},
ylabel={Pressure (p)},
axis background/.style={fill=white},
title style={font=\bfseries\color{mycolor1}},
title={Lax-Friedrichs Pressure Distribution},
xmajorgrids,
ymajorgrids,
legend style={legend cell align=left, align=left}
]
\addplot [color=green, line width=2.0pt]
  table[row sep=crcr]{%
0	0.999997562147529\\
0.0050251256281407	0.999997562147529\\
0.0100502512562814	0.999996021254339\\
0.0150753768844221	0.999996021254339\\
0.0201005025125628	0.999991544583892\\
0.0251256281407035	0.999991544583892\\
0.0301507537688442	0.9999822636553\\
0.0351758793969849	0.9999822636553\\
0.0402010050251256	0.999964029627012\\
0.0452261306532663	0.999964029627012\\
0.050251256281407	0.999929512449228\\
0.0552763819095477	0.999929512449228\\
0.0603015075376884	0.999866449552029\\
0.0653266331658292	0.999866449552029\\
0.0703517587939698	0.999755180279986\\
0.0753768844221105	0.999755180279986\\
0.0804020100502513	0.999565476732363\\
0.085427135678392	0.999565476732363\\
0.0904522613065327	0.999252794773278\\
0.0954773869346734	0.999252794773278\\
0.100502512562814	0.998754270286759\\
0.105527638190955	0.998754270286759\\
0.110552763819095	0.997985024715732\\
0.115577889447236	0.997985024715732\\
0.120603015075377	0.9968355589178\\
0.125628140703518	0.9968355589178\\
0.130653266331658	0.995171123208581\\
0.135678391959799	0.995171123208581\\
0.14070351758794	0.992833871780376\\
0.14572864321608	0.992833871780376\\
0.150753768844221	0.989648293992091\\
0.155778894472362	0.989648293992091\\
0.160804020100503	0.985429885267673\\
0.165829145728643	0.985429885267673\\
0.170854271356784	0.979996381866661\\
0.175879396984925	0.979996381866661\\
0.180904522613065	0.973180303438463\\
0.185929648241206	0.973180303438463\\
0.190954773869347	0.964841197356966\\
0.195979899497487	0.964841197356966\\
0.201005025125628	0.954875970021523\\
0.206030150753769	0.954875970021523\\
0.21105527638191	0.943226026603689\\
0.21608040201005	0.943226026603689\\
0.221105527638191	0.929880519789545\\
0.226130653266332	0.929880519789545\\
0.231155778894472	0.914875665375506\\
0.236180904522613	0.914875665375506\\
0.241206030150754	0.89829065481281\\
0.246231155778894	0.89829065481281\\
0.251256281407035	0.880241071899891\\
0.256281407035176	0.880241071899891\\
0.261306532663317	0.860870867364795\\
0.266331658291457	0.860870867364795\\
0.271356783919598	0.840343890577654\\
0.276381909547739	0.840343890577654\\
0.281407035175879	0.818835788865236\\
0.28643216080402	0.818835788865236\\
0.291457286432161	0.796526836226483\\
0.296482412060302	0.796526836226483\\
0.301507537688442	0.77359600566993\\
0.306532663316583	0.77359600566993\\
0.311557788944724	0.750216391499326\\
0.316582914572864	0.750216391499326\\
0.321608040201005	0.726551936335834\\
0.326633165829146	0.726551936335834\\
0.331658291457286	0.702755322878266\\
0.336683417085427	0.702755322878266\\
0.341708542713568	0.678966843357537\\
0.346733668341709	0.678966843357537\\
0.351758793969849	0.655314047692121\\
0.35678391959799	0.655314047692121\\
0.361809045226131	0.631911982023562\\
0.366834170854271	0.631911982023562\\
0.371859296482412	0.608863852254264\\
0.376884422110553	0.608863852254264\\
0.381909547738693	0.586261974802805\\
0.386934673366834	0.586261974802805\\
0.391959798994975	0.564188903964394\\
0.396984924623116	0.564188903964394\\
0.402010050251256	0.542718648967501\\
0.407035175879397	0.542718648967501\\
0.412060301507538	0.521917912429136\\
0.417085427135678	0.521917912429136\\
0.422110552763819	0.501847294717621\\
0.42713567839196	0.501847294717621\\
0.4321608040201	0.482562415561953\\
0.437185929648241	0.482562415561953\\
0.442211055276382	0.464114905275842\\
0.447236180904523	0.464114905275842\\
0.452261306532663	0.446553213697661\\
0.457286432160804	0.446553213697661\\
0.462311557788945	0.429923176396491\\
0.467336683417085	0.429923176396491\\
0.472361809045226	0.414268266742812\\
0.477386934673367	0.414268266742812\\
0.482412060301508	0.399629452372736\\
0.487437185929648	0.399629452372736\\
0.492462311557789	0.386044570631609\\
0.49748743718593	0.386044570631609\\
0.50251256281407	0.373547147287251\\
0.507537688442211	0.373547147287251\\
0.512562814070352	0.362164615499507\\
0.517587939698492	0.362164615499507\\
0.522613065326633	0.351915956957871\\
0.527638190954774	0.351915956957871\\
0.532663316582915	0.342808889347643\\
0.537688442211055	0.342808889347643\\
0.542713567839196	0.334836858806792\\
0.547738693467337	0.334836858806792\\
0.552763819095477	0.327976241296203\\
0.557788944723618	0.327976241296203\\
0.562814070351759	0.322184272358208\\
0.5678391959799	0.322184272358208\\
0.57286432160804	0.317398257426597\\
0.577889447236181	0.317398257426597\\
0.582914572864322	0.313536516891631\\
0.587939698492462	0.313536516891631\\
0.592964824120603	0.31050127504426\\
0.597989949748744	0.31050127504426\\
0.603015075376884	0.308183345880475\\
0.608040201005025	0.308183345880475\\
0.613065326633166	0.306468090593386\\
0.618090452261307	0.306468090593386\\
0.623115577889447	0.305241835514045\\
0.628140703517588	0.305241835514045\\
0.633165829145729	0.304397838031403\\
0.638190954773869	0.304397838031403\\
0.64321608040201	0.303841002508293\\
0.648241206030151	0.303841002508293\\
0.653266331658292	0.303490835330482\\
0.658291457286432	0.303490835330482\\
0.663316582914573	0.303282492261658\\
0.668341708542714	0.303282492261658\\
0.673366834170854	0.30316610409622\\
0.678391959798995	0.30316610409622\\
0.683417085427136	0.303104784727934\\
0.688442211055276	0.303104784727934\\
0.693467336683417	0.303071788982059\\
0.698492462311558	0.303071788982059\\
0.703517587939699	0.303047196260897\\
0.708542713567839	0.303047196260897\\
0.71356783919598	0.303014273299197\\
0.718592964824121	0.303014273299197\\
0.723618090452261	0.302955336635784\\
0.728643216080402	0.302955336635784\\
0.733668341708543	0.302846489759687\\
0.738693467336683	0.302846489759687\\
0.743718592964824	0.302650008548988\\
0.748743718592965	0.302650008548988\\
0.753768844221106	0.302302310755431\\
0.758793969849246	0.302302310755431\\
0.763819095477387	0.301694291414098\\
0.768844221105528	0.301694291414098\\
0.773869346733668	0.300639394189702\\
0.778894472361809	0.300639394189702\\
0.78391959798995	0.298823704798339\\
0.78894472361809	0.298823704798339\\
0.793969849246231	0.295733717388382\\
0.798994974874372	0.295733717388382\\
0.804020100502513	0.290567092827512\\
0.809045226130653	0.290567092827512\\
0.814070351758794	0.282162510234559\\
0.819095477386935	0.282162510234559\\
0.824120603015075	0.269054740678023\\
0.829145728643216	0.269054740678023\\
0.834170854271357	0.249862883285666\\
0.839195979899497	0.249862883285666\\
0.844221105527638	0.224219418319934\\
0.849246231155779	0.224219418319934\\
0.85427135678392	0.194002456049316\\
0.85929648241206	0.194002456049316\\
0.864321608040201	0.163662688613954\\
0.869346733668342	0.163662688613954\\
0.874371859296482	0.138319589432588\\
0.879396984924623	0.138319589432588\\
0.884422110552764	0.120734310869637\\
0.889447236180904	0.120734310869637\\
0.894472361809045	0.110335928173583\\
0.899497487437186	0.110335928173583\\
0.904522613065327	0.104870291326925\\
0.909547738693467	0.104870291326925\\
0.914572864321608	0.102209883488706\\
0.919597989949749	0.102209883488706\\
0.924623115577889	0.100975690838694\\
0.92964824120603	0.100975690838694\\
0.934673366834171	0.100421107617311\\
0.939698492462312	0.100421107617311\\
0.944723618090452	0.100177904608843\\
0.949748743718593	0.100177904608843\\
0.954773869346734	0.10007354276037\\
0.959798994974874	0.10007354276037\\
0.964824120603015	0.100029714790825\\
0.969849246231156	0.100029714790825\\
0.974874371859296	0.100011729969812\\
0.979899497487437	0.100011729969812\\
0.984924623115578	0.100004612266805\\
0.989949748743719	0.100004612266805\\
0.994974874371859	0.100002506498085\\
1	0.100002506498085\\
};
\addlegendentry{LF Solution}

\end{axis}
\end{tikzpicture}%}
    };
    
\end{tikzpicture}
\caption{Kết quả phương pháp Lax-Friedrichs tại t = 0.2}
\label{fig:lf_results}
\end{figure}
\subsubsection{Phương pháp local Lax-Friedrich}

Phương pháp local Lax-Friedrich (LLF) là phương pháp cải tiến cho Lax-Friedrich. Phương pháp này hoạt động bằng cách thay thế giá trị $\dfrac{\Delta t}{\Delta x}$ bằng một giá trị cục bộ $\alpha$. Thành phần được thêm vào này được gọi là vận tốc truyền sóng cực đại (Maximum Local Wave Speed):
\[
F_{i+1/2}^{LLF}(W_-, W_+) = \frac{1}{2} \left( f(W_-) + f(W_+) \right) - \dfrac{\alpha_{1/2}}{2} \left( W_+ - W_- \right).
\]
Hệ số $\alpha_{i+1/2}$ được tính bằng công thức
\[
\alpha_{i+1/2} = \max \left( |f'(W_i)|, |f'(W_{i+1})| \right)
\]
đối với phương trình vô hướng, còn đối với phương trình Euler thì ta có:
\[
\alpha_{i+1/2} = \max \left( |\lambda(A(W_i))|, |\lambda(A(W_{i+1}))| \right)
\]

\textbf{Ghi chú:}
Trong công thức tính bên trên, $A(W) = \frac{\partial f}{\partial W}$ là ma trận Jacobian của hệ phương trình. Các trị riêng $\lambda$ (eigenvalues) của ma trận này đại diện cho các vận tốc truyền sóng đặc trưng trong dòng chảy.
Đối với hệ phương trình Euler 1 chiều, ba trị riêng này lần lượt là $u-c$, $u$, và $u+c$ (trong đó $u$ là vận tốc dòng chảy và $c$ là vận tốc âm thanh). Do đó, để đảm bảo tính ổn định, hệ số $\alpha$ thực tế thường được tính bằng vận tốc sóng lớn nhất về mặt trị tuyệt đối:
\begin{equation*}
    \alpha_{i+1/2} = \max \left( |u_i| + c_i, |u_{i+1}| + c_{i+1} \right)
\end{equation*}

Dưới đây là kết quả áp dụng phương pháp Local Lax-Friedrich cho bài toán ống sốc Sod:

\begin{tikzpicture}
    % Define spacing and dimensions
    \def\plotwidth{7cm}
    \def\plotheight{5.5cm}
    \def\hspace{8cm}
    \def\vspace{6.5cm}
    
    % Top row
    \node[anchor=south west, fill=white] at (0,\vspace) {
        \resizebox{\plotwidth}{\plotheight}{% This file was created by matlab2tikz.
%
%The latest updates can be retrieved from
%  http://www.mathworks.com/matlabcentral/fileexchange/22022-matlab2tikz-matlab2tikz
%where you can also make suggestions and rate matlab2tikz.
%
\definecolor{mycolor1}{rgb}{0.12941,0.12941,0.12941}%
%
\begin{tikzpicture}

\begin{axis}[%
width=3.229in,
height=2.498in,
at={(0.542in,0.392in)},
scale only axis,
xmin=0,
xmax=1,
xlabel style={font=\color{mycolor1}},
xlabel={x},
ymin=0.1,
ymax=1,
ylabel style={font=\color{mycolor1}},
ylabel={$\text{Density (}\rho\text{)}$},
axis background/.style={fill=white},
title style={font=\bfseries\color{mycolor1}},
title={Local Lax-Friedrichs Density Distribution},
xmajorgrids,
ymajorgrids,
legend style={legend cell align=left, align=left}
]
\addplot [color=blue, line width=2.0pt]
  table[row sep=crcr]{%
0	1\\
0.0050251256281407	1\\
0.0100502512562814	1\\
0.0150753768844221	1\\
0.0201005025125628	0.999999999999999\\
0.0251256281407035	0.999999999999997\\
0.0301507537688442	0.999999999999989\\
0.0351758793969849	0.999999999999964\\
0.0402010050251256	0.999999999999888\\
0.0452261306532663	0.999999999999657\\
0.050251256281407	0.999999999998975\\
0.0552763819095477	0.999999999997004\\
0.0603015075376884	0.999999999991443\\
0.0653266331658292	0.999999999976125\\
0.0703517587939698	0.999999999934928\\
0.0753768844221105	0.999999999826737\\
0.0804020100502513	0.999999999549314\\
0.085427135678392	0.999999998854747\\
0.0904522613065327	0.999999997156958\\
0.0954773869346734	0.999999993105318\\
0.100502512562814	0.999999983666239\\
0.105527638190955	0.999999962200332\\
0.110552763819095	0.999999914551219\\
0.115577889447236	0.999999811321144\\
0.120603015075377	0.999999593069089\\
0.125628140703518	0.999999142812202\\
0.130653266331658	0.99999823654624\\
0.135678391959799	0.999996457114649\\
0.14070351758794	0.999993049365914\\
0.14572864321608	0.999986685255945\\
0.150753768844221	0.999975097284679\\
0.155778894472362	0.999954529618177\\
0.160804020100503	0.999918952442116\\
0.165829145728643	0.999858992495848\\
0.170854271356784	0.999760558763502\\
0.175879396984925	0.999603194059077\\
0.180904522613065	0.999358264246383\\
0.185929648241206	0.998987202373217\\
0.190954773869347	0.998440137466532\\
0.195979899497487	0.997655324500953\\
0.201005025125628	0.996559809383853\\
0.206030150753769	0.995071668540138\\
0.21105527638191	0.993103936743349\\
0.21608040201005	0.990570002393637\\
0.221105527638191	0.987389882067243\\
0.226130653266332	0.983496498790079\\
0.231155778894472	0.978840990504969\\
0.236180904522613	0.973396221490978\\
0.241206030150754	0.9671580276147\\
0.246231155778894	0.96014418461088\\
0.251256281407035	0.952391504430934\\
0.256281407035176	0.943951725761638\\
0.261306532663317	0.934886929482997\\
0.266331658291457	0.925265108408554\\
0.271356783919598	0.91515632535584\\
0.276381909547739	0.904629680366\\
0.281407035175879	0.893751130096939\\
0.28643216080402	0.882582083248085\\
0.291457286432161	0.871178634556589\\
0.296482412060302	0.85959128287754\\
0.301507537688442	0.847864989440668\\
0.306532663316583	0.836039456362968\\
0.311557788944724	0.824149533134098\\
0.316582914572864	0.81222568452705\\
0.321608040201005	0.80029447475747\\
0.326633165829146	0.788379039162617\\
0.331658291457286	0.776499526621133\\
0.336683417085427	0.764673504202449\\
0.341708542713568	0.752916320990961\\
0.346733668341709	0.741241431447049\\
0.351758793969849	0.729660680664244\\
0.35678391959799	0.718184554926768\\
0.361809045226131	0.706822401401257\\
0.366834170854271	0.695582620844885\\
0.371859296482412	0.684472837037361\\
0.376884422110553	0.673500046349147\\
0.381909547738693	0.662670750507383\\
0.386934673366834	0.65199107525281\\
0.391959798994975	0.641466877216275\\
0.396984924623116	0.631103840991312\\
0.402010050251256	0.620907568041382\\
0.407035175879397	0.610883658752917\\
0.412060301507538	0.601037788621546\\
0.417085427135678	0.591375779229376\\
0.422110552763819	0.581903664325378\\
0.42713567839196	0.572627750947056\\
0.4321608040201	0.563554675107921\\
0.437185929648241	0.554691451110757\\
0.442211055276382	0.546045513022107\\
0.447236180904523	0.537624746253808\\
0.452261306532663	0.529437506543961\\
0.457286432160804	0.521492622924039\\
0.462311557788945	0.513799380526776\\
0.467336683417085	0.506367478376793\\
0.472361809045226	0.499206956683551\\
0.477386934673367	0.492328087725128\\
0.482412060301508	0.485741224305027\\
0.487437185929648	0.479456600146437\\
0.492462311557789	0.47348407764299\\
0.49748743718593	0.467832840294842\\
0.50251256281407	0.462511030069879\\
0.507537688442211	0.457525333901853\\
0.512562814070352	0.452880528481975\\
0.517587939698492	0.448578998120242\\
0.522613065326633	0.444620246194288\\
0.527638190954774	0.441000425754481\\
0.532663316582915	0.437711918206844\\
0.537688442211055	0.434742989598446\\
0.542713567839196	0.432077551024855\\
0.547738693467337	0.429695042689764\\
0.552763819095477	0.427570450551401\\
0.557788944723618	0.425674451578859\\
0.562814070351759	0.423973670567253\\
0.5678391959799	0.422431020928558\\
0.57286432160804	0.421006096595256\\
0.577889447236181	0.419655584187555\\
0.582914572864322	0.418333674650165\\
0.587939698492462	0.416992470724394\\
0.592964824120603	0.415582408230758\\
0.597989949748744	0.414052731174269\\
0.603015075376884	0.412352078419027\\
0.608040201005025	0.410429248493948\\
0.613065326633166	0.408234205318168\\
0.618090452261307	0.405719369271635\\
0.623115577889447	0.402841205280994\\
0.628140703517588	0.399562075119549\\
0.633165829145729	0.395852270060315\\
0.638190954773869	0.39169208954316\\
0.64321608040201	0.387073789988982\\
0.648241206030151	0.382003203613411\\
0.653266331658292	0.376500826731702\\
0.658291457286432	0.370602204092318\\
0.663316582914573	0.364357489358146\\
0.668341708542714	0.357830136344351\\
0.673366834170854	0.351094761264962\\
0.678391959798995	0.344234300790533\\
0.683417085427136	0.337336661769808\\
0.688442211055276	0.330491105910547\\
0.693467336683417	0.323784630695033\\
0.698492462311558	0.317298595675064\\
0.703517587939699	0.311105805340701\\
0.708542713567839	0.305268203884933\\
0.71356783919598	0.29983527293273\\
0.718592964824121	0.294843159814459\\
0.723618090452261	0.290314508400915\\
0.728643216080402	0.286258921212626\\
0.733668341708543	0.282673951935917\\
0.738693467336683	0.279546510815825\\
0.743718592964824	0.276854559490986\\
0.748743718592965	0.274568974094158\\
0.753768844221106	0.272655463511551\\
0.758793969849246	0.271076441767082\\
0.763819095477387	0.269792768241435\\
0.768844221105528	0.268765285612988\\
0.773869346733668	0.267956101225267\\
0.778894472361809	0.267329569654073\\
0.78391959798995	0.26685293564161\\
0.78894472361809	0.26649657264903\\
0.793969849246231	0.266233671319262\\
0.798994974874372	0.266039023383904\\
0.804020100502513	0.26588605053043\\
0.809045226130653	0.265740132096625\\
0.814070351758794	0.265543742548805\\
0.819095477386935	0.265183652325222\\
0.824120603015075	0.264420149374002\\
0.829145728643216	0.262742939241247\\
0.834170854271357	0.259115416364568\\
0.839195979899497	0.251654731980178\\
0.844221105527638	0.237632147603155\\
0.849246231155779	0.214758725651435\\
0.85427135678392	0.184519029909397\\
0.85929648241206	0.155054435428736\\
0.864321608040201	0.136042954688121\\
0.869346733668342	0.128152002028822\\
0.874371859296482	0.125791229169379\\
0.879396984924623	0.12518970362446\\
0.884422110552764	0.125044858224337\\
0.889447236180904	0.125010562568674\\
0.894472361809045	0.125002482789629\\
0.899497487437186	0.125000582867147\\
0.904522613065327	0.12500013665775\\
0.909547738693467	0.12500003199132\\
0.914572864321608	0.125000007475209\\
0.919597989949749	0.12500000174276\\
0.924623115577889	0.125000000405205\\
0.92964824120603	0.125000000093909\\
0.934673366834171	0.12500000002168\\
0.939698492462312	0.125000000004983\\
0.944723618090452	0.125000000001139\\
0.949748743718593	0.125000000000259\\
0.954773869346734	0.125000000000058\\
0.959798994974874	0.125000000000013\\
0.964824120603015	0.125000000000003\\
0.969849246231156	0.125000000000001\\
0.974874371859296	0.125\\
0.979899497487437	0.125\\
0.984924623115578	0.125\\
0.989949748743719	0.125\\
0.994974874371859	0.125\\
1	0.125\\
};
\addlegendentry{LLF Solution}

\end{axis}
\end{tikzpicture}%}
    };
    
    \node[anchor=south west, fill=white] at (\hspace,\vspace) {
        \resizebox{\plotwidth}{\plotheight}{% This file was created by matlab2tikz.
%
%The latest updates can be retrieved from
%  http://www.mathworks.com/matlabcentral/fileexchange/22022-matlab2tikz-matlab2tikz
%where you can also make suggestions and rate matlab2tikz.
%
\definecolor{mycolor1}{rgb}{0.12941,0.12941,0.12941}%
%
\begin{tikzpicture}

\begin{axis}[%
width=3.229in,
height=2.498in,
at={(0.542in,0.392in)},
scale only axis,
xmin=0,
xmax=1,
xlabel style={font=\color{mycolor1}},
xlabel={x},
ymin=0,
ymax=1,
ylabel style={font=\color{mycolor1}},
ylabel={Velocity (u)},
axis background/.style={fill=white},
title style={font=\bfseries\color{mycolor1}},
title={Local Lax-Friedrichs Velocity Distribution},
xmajorgrids,
ymajorgrids,
legend style={legend cell align=left, align=left}
]
\addplot [color=red, line width=2.0pt]
  table[row sep=crcr]{%
0	0\\
0.0050251256281407	0\\
0.0100502512562814	9.75685235752145e-17\\
0.0150753768844221	3.21769486609396e-16\\
0.0201005025125628	1.1798257496155e-15\\
0.0251256281407035	4.00726354206742e-15\\
0.0301507537688442	1.31040585007194e-14\\
0.0351758793969849	4.20801238307517e-14\\
0.0402010050251256	1.32143601392406e-13\\
0.0452261306532663	4.05044395062861e-13\\
0.050251256281407	1.21252671629791e-12\\
0.0552763819095477	3.54516478701844e-12\\
0.0603015075376884	1.01250366810572e-11\\
0.0653266331658292	2.82489591795703e-11\\
0.0703517587939698	7.699413493165e-11\\
0.0753768844221105	2.05007108749359e-10\\
0.0804020100502513	5.33259299640377e-10\\
0.085427135678392	1.35508172934608e-09\\
0.0904522613065327	3.36393263392178e-09\\
0.0954773869346734	8.15789834696141e-09\\
0.100502512562814	1.93263662190758e-08\\
0.105527638190955	4.47251706848823e-08\\
0.110552763819095	1.01104360561949e-07\\
0.115577889447236	2.23247833282821e-07\\
0.120603015075377	4.81487151863619e-07\\
0.125628140703518	1.01423832241445e-06\\
0.130653266331658	2.08654688493209e-06\\
0.135678391959799	4.19199982554654e-06\\
0.14070351758794	8.22410745941132e-06\\
0.14572864321608	1.57542446600663e-05\\
0.150753768844221	2.94653977588645e-05\\
0.155778894472362	5.38016816190077e-05\\
0.160804020100503	9.58981609661811e-05\\
0.165829145728643	0.000166846919967129\\
0.170854271356784	0.000283324922300638\\
0.175879396984925	0.000469548794899554\\
0.180904522613065	0.000759427417747498\\
0.185929648241206	0.00119866102703121\\
0.190954773869347	0.00184640639652912\\
0.195979899497487	0.00277602960336992\\
0.201005025125628	0.00407445086209352\\
0.206030150753769	0.00583969658058496\\
0.21105527638191	0.00817653208850657\\
0.21608040201005	0.0111904248074598\\
0.221105527638191	0.0149804949433606\\
0.226130653266332	0.0196324240177806\\
0.231155778894472	0.0252123938122629\\
0.236180904522613	0.0317629649185391\\
0.241206030150754	0.0393014146016345\\
0.246231155778894	0.0478205597761447\\
0.251256281407035	0.0572916440274085\\
0.256281407035176	0.0676685852678408\\
0.261306532663317	0.0788928058384341\\
0.266331658291457	0.0908979679942946\\
0.271356783919598	0.103614139662612\\
0.276381909547739	0.116971138889329\\
0.281407035175879	0.130900994207244\\
0.28643216080402	0.145339587029168\\
0.291457286432161	0.160227610851083\\
0.296482412060302	0.175511003919112\\
0.301507537688442	0.191141004101316\\
0.306532663316583	0.207073951776289\\
0.311557788944724	0.223270938928866\\
0.316582914572864	0.239697376342693\\
0.321608040201005	0.256322528587388\\
0.326633165829146	0.273119049176388\\
0.331658291457286	0.290062535505348\\
0.336683417085427	0.307131114205417\\
0.341708542713568	0.324305061502743\\
0.346733668341709	0.341566459286973\\
0.351758793969849	0.358898885222535\\
0.35678391959799	0.376287133902307\\
0.361809045226131	0.393716965390473\\
0.366834170854271	0.411174877281525\\
0.371859296482412	0.428647896446621\\
0.376884422110553	0.446123386835716\\
0.381909547738693	0.463588869984379\\
0.386934673366834	0.481031855196838\\
0.391959798994975	0.498439676720742\\
0.396984924623116	0.515799335587681\\
0.402010050251256	0.533097344170021\\
0.407035175879397	0.550319571910378\\
0.412060301507538	0.56745109113246\\
0.417085427135678	0.584476022363851\\
0.422110552763819	0.601377379219884\\
0.42713567839196	0.618136913644634\\
0.4321608040201	0.634734963214182\\
0.437185929648241	0.651150303313811\\
0.442211055276382	0.667360008337082\\
0.447236180904523	0.683339327645464\\
0.452261306532663	0.699061583881342\\
0.457286432160804	0.714498103327091\\
0.462311557788945	0.729618190288981\\
0.467336683417085	0.744389159838051\\
0.472361809045226	0.758776445460569\\
0.477386934673367	0.772743799957182\\
0.482412060301508	0.786253608863789\\
0.487437185929648	0.799267335210293\\
0.492462311557789	0.811746111948821\\
0.49748743718593	0.823651493192855\\
0.50251256281407	0.834946366901865\\
0.507537688442211	0.845596019440758\\
0.512562814070352	0.855569326600176\\
0.517587939698492	0.864840026917424\\
0.522613065326633	0.873388013094332\\
0.527638190954774	0.881200558518037\\
0.532663316582915	0.888273381706049\\
0.537688442211055	0.894611445637447\\
0.542713567839196	0.900229394749933\\
0.547738693467337	0.905151551932894\\
0.552763819095477	0.909411430988049\\
0.557788944723618	0.913050763864189\\
0.562814070351759	0.916118090921565\\
0.5678391959799	0.918667009129846\\
0.57286432160804	0.920754209707608\\
0.577889447236181	0.92243745698256\\
0.582914572864322	0.923773660883125\\
0.587939698492462	0.924817176918371\\
0.592964824120603	0.925618433769996\\
0.597989949748744	0.926222946158469\\
0.603015075376884	0.926670726739878\\
0.608040201005025	0.926996071949621\\
0.613065326633166	0.92722766749561\\
0.618090452261307	0.927388941724959\\
0.623115577889447	0.9274985891368\\
0.628140703517588	0.927571189961866\\
0.633165829145729	0.927617862113254\\
0.638190954773869	0.927646895877924\\
0.64321608040201	0.927664336772806\\
0.648241206030151	0.927674496014739\\
0.653266331658292	0.927680379794195\\
0.658291457286432	0.927684037437565\\
0.663316582914573	0.927686834566497\\
0.668341708542714	0.927689660843553\\
0.673366834170854	0.927693083343571\\
0.678391959798995	0.92769745657383\\
0.683417085427136	0.92770299920986\\
0.688442211055276	0.927709846155864\\
0.693467336683417	0.927718082909038\\
0.698492462311558	0.927727767629344\\
0.703517587939699	0.92773894492163\\
0.708542713567839	0.927751654184835\\
0.71356783919598	0.92776593448206\\
0.718592964824121	0.927781827212588\\
0.723618090452261	0.927799377383456\\
0.728643216080402	0.927818633938493\\
0.733668341708543	0.927839649361042\\
0.738693467336683	0.92786247857514\\
0.743718592964824	0.927887176970011\\
0.748743718592965	0.927913797068943\\
0.753768844221106	0.927942382764299\\
0.758793969849246	0.927972958719226\\
0.763819095477387	0.928005509513935\\
0.768844221105528	0.928039936120443\\
0.773869346733668	0.928075961072688\\
0.778894472361809	0.928112916149496\\
0.78391959798995	0.928149259686866\\
0.78894472361809	0.928181471097922\\
0.793969849246231	0.928201512249665\\
0.798994974874372	0.928190996670724\\
0.804020100502513	0.928107816293345\\
0.809045226130653	0.927855675317821\\
0.814070351758794	0.927215002092284\\
0.819095477386935	0.925687890361029\\
0.824120603015075	0.922158502206075\\
0.829145728643216	0.914185631693572\\
0.834170854271357	0.896676688224371\\
0.839195979899497	0.859903519536217\\
0.844221105527638	0.787942642043329\\
0.849246231155779	0.661364343002731\\
0.85427135678392	0.472096246027225\\
0.85929648241206	0.255908552279432\\
0.864321608040201	0.0965342738062098\\
0.869346733668342	0.027243046761056\\
0.874371859296482	0.00675255830512803\\
0.879396984924623	0.00161011011153123\\
0.884422110552764	0.000380054310940096\\
0.889447236180904	8.94436960458631e-05\\
0.894472361809045	2.10213128754121e-05\\
0.899497487437186	4.93484870000321e-06\\
0.904522613065327	1.15700322337851e-06\\
0.909547738693467	2.70851642617963e-07\\
0.914572864321608	6.32881489862829e-08\\
0.919597989949749	1.47549134339687e-08\\
0.924623115577889	3.4306332785133e-09\\
0.92964824120603	7.95068990551865e-10\\
0.934673366834171	1.8355514769541e-10\\
0.939698492462312	4.21858871908449e-11\\
0.944723618090452	9.64455208133057e-12\\
0.949748743718593	2.19150286498935e-12\\
0.954773869346734	4.94444957285347e-13\\
0.959798994974874	1.106086435552e-13\\
0.964824120603015	2.45175204564859e-14\\
0.969849246231156	5.37169533284924e-15\\
0.974874371859296	1.10605688142565e-15\\
0.979899497487437	1.9513734897128e-16\\
0.984924623115578	0\\
0.989949748743719	0\\
0.994974874371859	0\\
1	0\\
};
\addlegendentry{LLF Solution}

\end{axis}
\end{tikzpicture}%}
    };
    
    % Bottom row
    \node[anchor=south west, fill=white] at (0,0) {
        \resizebox{\plotwidth}{\plotheight}{% This file was created by matlab2tikz.
%
%The latest updates can be retrieved from
%  http://www.mathworks.com/matlabcentral/fileexchange/22022-matlab2tikz-matlab2tikz
%where you can also make suggestions and rate matlab2tikz.
%
\definecolor{mycolor1}{rgb}{0.12941,0.12941,0.12941}%
%
\begin{tikzpicture}

\begin{axis}[%
width=3.229in,
height=2.498in,
at={(0.542in,0.392in)},
scale only axis,
xmin=0,
xmax=1,
xlabel style={font=\color{mycolor1}},
xlabel={x},
ymin=0.1,
ymax=1,
ylabel style={font=\color{mycolor1}},
ylabel={Pressure (p)},
axis background/.style={fill=white},
title style={font=\bfseries\color{mycolor1}},
title={Local Lax-Friedrichs Pressure Distribution},
xmajorgrids,
ymajorgrids,
legend style={legend cell align=left, align=left}
]
\addplot [color=green, line width=2.0pt]
  table[row sep=crcr]{%
0	1\\
0.0050251256281407	1\\
0.0100502512562814	1\\
0.0150753768844221	0.999999999999999\\
0.0201005025125628	0.999999999999999\\
0.0251256281407035	0.999999999999995\\
0.0301507537688442	0.999999999999984\\
0.0351758793969849	0.99999999999995\\
0.0402010050251256	0.999999999999844\\
0.0452261306532663	0.999999999999521\\
0.050251256281407	0.999999999998565\\
0.0552763819095477	0.999999999995805\\
0.0603015075376884	0.99999999998802\\
0.0653266331658292	0.999999999966576\\
0.0703517587939698	0.9999999999089\\
0.0753768844221105	0.999999999757432\\
0.0804020100502513	0.999999999369039\\
0.085427135678392	0.999999998396646\\
0.0904522613065327	0.999999996019741\\
0.0954773869346734	0.999999990347445\\
0.100502512562814	0.999999977132735\\
0.105527638190955	0.999999947080465\\
0.110552763819095	0.999999880371713\\
0.115577889447236	0.999999735849628\\
0.120603015075377	0.999999430296844\\
0.125628140703518	0.999998799937592\\
0.130653266331658	0.999997531166808\\
0.135678391959799	0.999995039968556\\
0.14070351758794	0.999990269142111\\
0.14572864321608	0.999981359463809\\
0.150753768844221	0.999965136554322\\
0.155778894472362	0.999936342609736\\
0.160804020100503	0.999886536928223\\
0.165829145728643	0.999802599754185\\
0.170854271356784	0.999664810862783\\
0.175879396984925	0.99944454763672\\
0.180904522613065	0.999101762236947\\
0.185929648241206	0.998582547324146\\
0.190954773869347	0.997817259753941\\
0.195979899497487	0.996719795213772\\
0.201005025125628	0.995188631065307\\
0.206030150753769	0.993110118618918\\
0.21105527638191	0.990364180029977\\
0.21608040201005	0.98683208188563\\
0.221105527638191	0.982405426857108\\
0.226130653266332	0.976995091182383\\
0.231155778894472	0.970538699153973\\
0.236180904522613	0.963005446209142\\
0.241206030150754	0.954397611014871\\
0.246231155778894	0.944748768888868\\
0.251256281407035	0.934119323772322\\
0.256281407035176	0.922590349585548\\
0.261306532663317	0.910256815276909\\
0.266331658291457	0.897221108572211\\
0.271356783919598	0.883587479633745\\
0.276381909547739	0.86945770989858\\
0.281407035175879	0.854928051134345\\
0.28643216080402	0.840087307989918\\
0.291457286432161	0.825015851473498\\
0.296482412060302	0.80978532987429\\
0.301507537688442	0.794458862547214\\
0.306532663316583	0.779091539643992\\
0.311557788944724	0.763731093024449\\
0.316582914572864	0.748418642239244\\
0.321608040201005	0.733189451223951\\
0.326633165829146	0.718073655542149\\
0.331658291457286	0.703096937423959\\
0.336683417085427	0.688281137757593\\
0.341708542713568	0.673644801944801\\
0.346733668341709	0.659203661295172\\
0.351758793969849	0.644971054337295\\
0.35678391959799	0.630958293760535\\
0.361809045226131	0.61717498516588\\
0.366834170854271	0.603629303740139\\
0.371859296482412	0.59032823460349\\
0.376884422110553	0.5772777820638\\
0.381909547738693	0.564483152434132\\
0.386934673366834	0.551948914486713\\
0.391959798994975	0.539679141056546\\
0.396984924623116	0.527677534783879\\
0.402010050251256	0.51594754049973\\
0.407035175879397	0.504492446308812\\
0.412060301507538	0.493315475001941\\
0.417085427135678	0.482419867025144\\
0.422110552763819	0.47180895583393\\
0.42713567839196	0.461486236056407\\
0.4321608040201	0.451455424466223\\
0.437185929648241	0.441720513314724\\
0.442211055276382	0.432285815082076\\
0.447236180904523	0.423155997174116\\
0.452261306532663	0.414336104515451\\
0.457286432160804	0.405831567378972\\
0.462311557788945	0.397648191168874\\
0.467336683417085	0.389792124277051\\
0.472361809045226	0.382269799621998\\
0.477386934673367	0.375087845142267\\
0.482412060301508	0.368252958470378\\
0.487437185929648	0.36177174140407\\
0.492462311557789	0.355650490788818\\
0.49748743718593	0.349894944203517\\
0.50251256281407	0.344509981552046\\
0.507537688442211	0.339499287392773\\
0.512562814070352	0.3348649835489\\
0.517587939698492	0.330607247012783\\
0.522613065326633	0.326723933925494\\
0.527638190954774	0.323210235752788\\
0.532663316582915	0.320058397730533\\
0.537688442211055	0.317257531131874\\
0.542713567839196	0.314793548901085\\
0.547738693467337	0.312649248025559\\
0.552763819095477	0.310804551607676\\
0.557788944723618	0.309236909696744\\
0.562814070351759	0.307921842146168\\
0.5678391959799	0.306833591337032\\
0.57286432160804	0.305945840073047\\
0.577889447236181	0.305232442518274\\
0.582914572864322	0.304668115092001\\
0.587939698492462	0.304229039931949\\
0.592964824120603	0.303893344817348\\
0.597989949748744	0.303641438261034\\
0.603015075376884	0.303456194316495\\
0.608040201005025	0.30332299608663\\
0.613065326633166	0.303229658132069\\
0.618090452261307	0.303166254972278\\
0.623115577889447	0.303124885539525\\
0.628140703517588	0.30309940237405\\
0.633165829145729	0.303085130554335\\
0.638190954773869	0.30307859599265\\
0.64321608040201	0.303077276851782\\
0.648241206030151	0.30307938627591\\
0.653266331658292	0.30308368990863\\
0.658291457286432	0.303089358051198\\
0.663316582914573	0.30309584983207\\
0.668341708542714	0.303102825304094\\
0.673366834170854	0.303110080765387\\
0.678391959798995	0.303117502593018\\
0.683417085427136	0.303125035272834\\
0.688442211055276	0.303132659923034\\
0.693467336683417	0.303140380303383\\
0.698492462311558	0.303148213978981\\
0.703517587939699	0.303156186908068\\
0.708542713567839	0.303164330219445\\
0.71356783919598	0.303172678331115\\
0.718592964824121	0.30318126784632\\
0.723618090452261	0.303190136862228\\
0.728643216080402	0.303199324457583\\
0.733668341708543	0.303208870204094\\
0.738693467336683	0.303218813581863\\
0.743718592964824	0.303229193171932\\
0.748743718592965	0.303240045431348\\
0.753768844221106	0.303251402673652\\
0.758793969849246	0.303263289443594\\
0.763819095477387	0.303275715464231\\
0.768844221105528	0.303288660986822\\
0.773869346733668	0.303302044925829\\
0.778894472361809	0.303315653544906\\
0.78391959798995	0.303328978332615\\
0.78894472361809	0.303340844676745\\
0.793969849246231	0.303348559140564\\
0.798994974874372	0.303345951114428\\
0.804020100502513	0.303318882660083\\
0.809045226130653	0.303235027610168\\
0.814070351758794	0.303020740365604\\
0.819095477386935	0.302509384549342\\
0.824120603015075	0.301329429112999\\
0.829145728643216	0.298677685690245\\
0.834170854271357	0.292924388239093\\
0.839195979899497	0.281152931395329\\
0.844221105527638	0.259303717948389\\
0.849246231155779	0.224442568522252\\
0.85427135678392	0.179901299599167\\
0.85929648241206	0.138490169954664\\
0.864321608040201	0.113359452753888\\
0.869346733668342	0.103646204222034\\
0.874371859296482	0.100895727645497\\
0.879396984924623	0.100213133634763\\
0.884422110552764	0.10005028398727\\
0.889447236180904	0.100011832704761\\
0.894472361809045	0.100002780881257\\
0.899497487437186	0.100000652820385\\
0.904522613065327	0.100000153057209\\
0.909547738693467	0.100000035830308\\
0.914572864321608	0.100000008372235\\
0.919597989949749	0.100000001951892\\
0.924623115577889	0.10000000045383\\
0.92964824120603	0.100000000105178\\
0.934673366834171	0.100000000024282\\
0.939698492462312	0.100000000005581\\
0.944723618090452	0.100000000001276\\
0.949748743718593	0.10000000000029\\
0.954773869346734	0.100000000000065\\
0.959798994974874	0.100000000000015\\
0.964824120603015	0.100000000000003\\
0.969849246231156	0.100000000000001\\
0.974874371859296	0.1\\
0.979899497487437	0.1\\
0.984924623115578	0.1\\
0.989949748743719	0.1\\
0.994974874371859	0.1\\
1	0.1\\
};
\addlegendentry{LLF Solution}

\end{axis}
\end{tikzpicture}%}
    };
    
\end{tikzpicture}

\subsection{Phương pháp bậc 2 cho không gian - Piecewise linear construction}
Phương pháp LLF ở trên cho ta sai số bậc một theo không gian, vì nó dựa trên việc tái cấu trúc từng đoạn bằng hằng số. Cụ thể hơn, trong công thức
\begin{equation*}
\begin{aligned}
    F_{i+1/2} &= \mathcal{F}(W_{i+1/2}^-, W_{i+1/2}^+) \\
    &= \frac{1}{2} \left( f(W_{i+1/2}^-) + f(W_{i+1/2}^+) \right) - \frac{\alpha_{i+1/2}}{2} (W_{i+1/2}^+ - W_{i+1/2}^-),
\end{aligned}
\end{equation*}
đối với LLF thì ta có $W_{i+1/2}^- = W_i, W_{i+1/2}^+ = W_{i+1}$.

Để tăng độ chính xác lên bậc 2 theo không gian, thay vì để $W_{i+1/2}^-$ và $W_{i+1/2}^+$ như trên, ta sẽ dùng các công thức sau:
\[
W_{i+1/2}^- = W_i + W_i' \frac{\Delta x}{2}, \quad W_{i+1/2}^+ = W_{i+1} - W_{i+1}' \frac{\Delta x}{2}.
\]
Phương pháp này được gọi là piecewise linear construction.
Trong công thức, $W_i'$ là đạo hàm tại $x_i$, tuy nhiên có thể được thay thế bởi một công thức tính toán phù hợp (do đạo hàm rất khó tính). Cách đơn giản nhất là dùng hàm minmod với hai tham số:
\[
W'_i = \frac{1}{\Delta x}\text{minmod}(W_i - W_{i-1}, W_{i+1} - W_i).
\]
Với hàm minmod được cho bởi
\begin{equation*}
    \begin{split}
    \text{minmod(a, b)} &= \begin{cases}
        a & \text{if } ab > 0, |a| \le |b|, \\
        b & \text{if } ab > 0, |b| < |a|, \\
        0 & \text{if } ab \le 0.
    \end{cases} \\
    % & = \frac{1}{2}(\sgn(a) + \sgn(b)) \min(|a|, |b|).
\end{split}
\end{equation*}
hoặc công thức rút gọn
\[
\text{minmod(a, b)} = \dfrac{1}{2}(\text{sgn}(a) + \text{sgn}(b)) \min(|a|, |b|).
\]

Bên cạnh công thức minmod cho hai tham số như trên thì ta cũng có thể dùng công thức minmod cho ba tham số:
\[
W'_i = \text{minmod3}\left(\frac{\theta(W_{i+1} - W_i)}{\Delta x}, \frac{(W_{i+1} - W_{i-1})}{2\Delta x}, \frac{\theta(W_i - W_{i-1})}{\Delta x}\right).
\]
và công thức minmod3 được cho bởi
\begin{equation*}
    \text{minmod3}(a, b, c) = \begin{cases}
        \min(|a|, |b|, |c|)\text{sign}(a) & \text{nếu } a, b, c \text{ cùng dấu}, \\
        0 & \text{nếu } a, b, c \text{ khác dấu}.
    \end{cases}
\end{equation*}

Dưới đây là kết quả áp dụng phương pháp bậc 2 không gian cho bài toán ống sốc Sod:

\input{figure/spatial2_results}

\subsection{Phương pháp bậc 2 theo thời gian - Công thức Heun}

\subsubsection{Biến đổi phương trình về dạng ODE}
Xét phương trình Euler 1 chiều:
\begin{equation*}
    w_t + f(w)_x = 0
\end{equation*}
trong đó: 
\[
w = \begin{pmatrix}
    \rho \\ \rho u \\ E
\end{pmatrix}, \quad f(w) = \begin{pmatrix}
    \rho u \\ \rho u^2 + p \\ (E + p)u
\end{pmatrix}.
\]
Trước tiên, ta lấy tích phân phương trình ban đầu trên đoạn $[x_{i-\frac{1}{2}}, x_{i+\frac{1}{2}}]$:
\[
\int_{x_{i-1/2}}^{x_{i+1/2}} \left( \frac{\partial w}{\partial t} + \frac{\partial f(w)}{\partial x} \right) dx = 0, \quad \text{hay} \quad \int_{x_{i-1/2}}^{x_{i+1/2}}  \frac{\partial w}{\partial t} dx + \int_{x_{i-1/2}}^{x_{i+1/2}} \frac{\partial f(w)}{\partial x}  dx = 0
\]
Áp dụng quy tắc Leibniz cho số hạng đầu tiên và định lý cơ bản của giải tích cho số hạng thứ hai:
\begin{align*}
        \int_{x_{i-1/2}}^{x_{i+1/2}}  \frac{\partial w}{\partial t} dx &= \frac{d}{dt} \int_{x_{i-1/2}}^{x_{i+1/2}} w(x,t) \, dx, \\
        \int_{x_{i-1/2}}^{x_{i+1/2}} \frac{\partial f(w)}{\partial x}  dx &= \left[ f(w(x_{i+1/2}, t)) - f(w(x_{i-1/2}, t)) \right].
\end{align*}
Suy ra:
\[
\frac{d}{dt} \int_{x_{i-1/2}}^{x_{i+1/2}} w(x,t) \, dx + \left[ f(w(x_{i+1/2}, t)) - f(w(x_{i-1/2}, t)) \right] = 0.
\]
Chia hai vế của phương trình trên cho $\Delta x$:
\[
\frac{1}{\Delta x}\frac{d}{dt} \int_{x_{i-1/2}}^{x_{i+1/2}} w(x,t) \, dx + \frac{1}{\Delta x}\left[ f(w(x_{i+1/2}, t)) - f(w(x_{i-1/2}, t)) \right] = 0.
\]
Đặt giá trị trung bình trên ô (cell average) và các thông lượng số (numerical flux):
\[
\langle w \rangle_i = \frac{1}{\Delta x}\int_{x_{i-1/2}}^{x_{i+1/2}} w(x,t) \, dx, \quad F_{i+1/2} \approx f(w(x_{i+1/2}, t)), \quad F_{i-1/2} \approx f(w(x_{i-1/2}, t)).
\]
Cuối cùng ta thu được phương trình:
\[
\frac{d\langle w \rangle_i}{dt} = - \frac{F_{i+1/2} - F_{i-1/2}}{\Delta x} =: \mathcal{F}_i(w).
\]
Hệ phương trình thu được sau cùng là hệ ODE và có thể được tính xấp xỉ nghiệm bằng cách dùng phương pháp số cho phương trình vi phân. Trong bài này, ta sẽ dùng công thức Heun để đạt được sai số bậc 2 về mặt thời gian.

\subsubsection{Áp dụng công thức Heun}
Công thức lặp tổng quát của phương pháp Heun để tìm nghiệm xấp xỉ $y_{i+1}$ tại bước $x_{i+1} = x_i + h$ (với $h$ là bước nhảy thời gian) được xác định như sau:
\begin{equation*}
    y_{i+1} = y_i + \frac{h}{2} \left[ f(x_i, y_i) + f\Big(x_i + h, \, y_i + h f(x_i, y_i)\Big)\right]
\end{equation*}
Trong đó, công thức có thể được diễn giải qua hai giai đoạn (Dự báo - Hiệu chỉnh):
\begin{align*}
    \text{Dự báo (Predictor):} & \quad \tilde{y}_{i+1} = y_i + h f(x_i, y_i) \\
    \text{Hiệu chỉnh (Corrector):} & \quad y_{i+1} = y_i + \frac{h}{2} \left[ f(x_i, y_i) + f(x_{i+1}, \tilde{y}_{i+1}) \right]
\end{align*}

\noindent Để áp dụng công thức này vào bài toán, trước tiên ta sẽ "lắp ghép" các vector biến bảo toàn $W_i (\approx \langle w \rangle_i)$ và vector tốc độ biến thiên $\mathcal{F}_i$ thành các Vector trạng thái toàn cục $W$ và $\mathcal{F}$, khi đó ta có hệ:
\[
\dfrac{dW}{dt} = \mathcal{F}(W).
\]
Áp dụng công thức Heun với bước thời gian $\Delta t$, ta có công thức cập nhật nghiệm:
\[
W^{n+1} = W^n + \frac{\Delta t}{2} \left[ \mathcal{F}(W^n) + \mathcal{F}\Big(W^n + \Delta t \mathcal{F}(W^n)\Big)\right].
\]
Cụ thể hơn, thuật toán được thực hiện qua 4 bước:
\begin{enumerate}
    \item Bước 1 (Tính độ dốc ban đầu):
    Tính toán vector thay đổi dựa trên trạng thái hiện tại $W^n$:
    \begin{equation*}
        K_1 = \mathcal{F}(W^n)
    \end{equation*}

    \item Bước 2 (Dự báo - Predictor):
    Ước lượng nghiệm trung gian $W^*$ (nghiệm dự báo) bằng phương pháp Euler thuận:
    \begin{equation*}
        W^* = W^n + \Delta t K_1
    \end{equation*}

    \item Bước 3 (Tính độ dốc tại điểm dự báo):
    Tính toán vector thay đổi dựa trên trạng thái dự báo $W^*$:
    \begin{equation*}
        K_2 = \mathcal{F}(W^*)
    \end{equation*}

    \item Bước 4 (Hiệu chỉnh - Corrector):
    Cập nhật nghiệm chính thức tại bước thời gian mới $n+1$ bằng trung bình cộng của hai độ dốc:
    \begin{equation*}
        W^{n+1} = W^n + \frac{\Delta t}{2} (K_1 + K_2)
    \end{equation*}
\end{enumerate}

Dưới đây là kết quả áp dụng phương pháp bậc 2 thời gian cho bài toán ống sốc Sod:

\begin{figure}[H]
\centering
\begin{tikzpicture}
    % Define spacing and dimensions
    \def\plotwidth{7cm}
    \def\plotheight{5.5cm}
    \def\hspace{8cm}
    \def\vspace{6.5cm}
    
    % Top row
    \node[anchor=south west, fill=white] at (0,\vspace) {
        \resizebox{\plotwidth}{\plotheight}{% This file was created by matlab2tikz.
%
%The latest updates can be retrieved from
%  http://www.mathworks.com/matlabcentral/fileexchange/22022-matlab2tikz-matlab2tikz
%where you can also make suggestions and rate matlab2tikz.
%
\definecolor{mycolor1}{rgb}{0.12941,0.12941,0.12941}%
%
\begin{tikzpicture}

\begin{axis}[%
width=3.229in,
height=2.498in,
at={(0.542in,0.392in)},
scale only axis,
xmin=0,
xmax=1,
xlabel style={font=\color{mycolor1}},
xlabel={x},
ymin=0.1,
ymax=1,
ylabel style={font=\color{mycolor1}},
ylabel={$\text{Density (}\rho\text{)}$},
axis background/.style={fill=white},
title style={font=\bfseries\color{mycolor1}},
title={Second-Order Temporal Density Distribution},
xmajorgrids,
ymajorgrids,
legend style={legend cell align=left, align=left}
]
\addplot [color=blue, line width=2.0pt]
  table[row sep=crcr]{%
0	0.999999999995799\\
0.0050251256281407	0.999999999995799\\
0.0100502512562814	0.999999999990777\\
0.0150753768844221	0.999999999979982\\
0.0201005025125628	0.999999999957054\\
0.0251256281407035	0.999999999908937\\
0.0301507537688442	0.99999999980917\\
0.0351758793969849	0.999999999604815\\
0.0402010050251256	0.999999999191343\\
0.0452261306532663	0.999999998365062\\
0.050251256281407	0.999999996734326\\
0.0552763819095477	0.999999993556247\\
0.0603015075376884	0.999999987440892\\
0.0653266331658292	0.999999975823653\\
0.0703517587939698	0.999999954038781\\
0.0753768844221105	0.99999991371861\\
0.0804020100502513	0.999999840072914\\
0.085427135678392	0.999999707342588\\
0.0904522613065327	0.999999471333651\\
0.0954773869346734	0.999999057374586\\
0.100502512562814	0.999998341248583\\
0.105527638190955	0.999997119577965\\
0.110552763819095	0.999995064738371\\
0.115577889447236	0.999991657644774\\
0.120603015075377	0.999986089732543\\
0.125628140703518	0.999977123310024\\
0.130653266331658	0.999962897490841\\
0.135678391959799	0.999940665626574\\
0.14070351758794	0.999906450281084\\
0.14572864321608	0.999854604256076\\
0.150753768844221	0.999777272061218\\
0.155778894472362	0.999663756529153\\
0.160804020100503	0.999499810642339\\
0.165829145728643	0.999266894903538\\
0.170854271356784	0.998941464309304\\
0.175879396984925	0.998494373074725\\
0.180904522613065	0.997890504842732\\
0.185929648241206	0.997088744996133\\
0.190954773869347	0.99604240348671\\
0.195979899497487	0.994700166544245\\
0.201005025125628	0.993007602753701\\
0.206030150753769	0.990909177880017\\
0.21105527638191	0.988350654143073\\
0.21608040201005	0.985281678670512\\
0.221105527638191	0.981658318928767\\
0.226130653266332	0.977445292748851\\
0.231155778894472	0.972617671892372\\
0.236180904522613	0.967161905731206\\
0.241206030150754	0.961076101534137\\
0.246231155778894	0.954369591374193\\
0.251256281407035	0.94706189485757\\
0.256281407035176	0.939181239249279\\
0.261306532663317	0.930762819483413\\
0.266331658291457	0.921846972815771\\
0.271356783919598	0.912477414385503\\
0.276381909547739	0.902699640488368\\
0.281407035175879	0.892559564892268\\
0.28643216080402	0.882102416553349\\
0.291457286432161	0.87137189819668\\
0.296482412060302	0.86040958543772\\
0.301507537688442	0.849254534704854\\
0.306532663316583	0.837943063512259\\
0.311557788944724	0.826508666688192\\
0.316582914572864	0.81498203519884\\
0.321608040201005	0.803391148790238\\
0.326633165829146	0.791761418774326\\
0.331658291457286	0.780115862250995\\
0.336683417085427	0.768475293515961\\
0.341708542713568	0.756858522190696\\
0.346733668341709	0.745282550694184\\
0.351758793969849	0.733762766102923\\
0.35678391959799	0.722313123298875\\
0.361809045226131	0.710946317680831\\
0.366834170854271	0.699673946706523\\
0.371859296482412	0.688506660223949\\
0.376884422110553	0.677454300011097\\
0.381909547738693	0.666526029229943\\
0.386934673366834	0.65573045265705\\
0.391959798994975	0.645075728612159\\
0.396984924623116	0.634569673491089\\
0.402010050251256	0.624219859735801\\
0.407035175879397	0.614033707951847\\
0.412060301507538	0.604018573715857\\
0.417085427135678	0.594181829403064\\
0.422110552763819	0.58453094110378\\
0.42713567839196	0.575073540382269\\
0.4321608040201	0.565817490254098\\
0.437185929648241	0.556770944310728\\
0.442211055276382	0.547942397396069\\
0.447236180904523	0.539340725635365\\
0.452261306532663	0.530975212934989\\
0.457286432160804	0.522855560325515\\
0.462311557788945	0.514991873738397\\
0.467336683417085	0.507394625038119\\
0.472361809045226	0.50007458045388\\
0.477386934673367	0.493042690077499\\
0.482412060301508	0.486309931962915\\
0.487437185929648	0.479887104756607\\
0.492462311557789	0.473784563909357\\
0.49748743718593	0.468011898569207\\
0.50251256281407	0.462577549394658\\
0.507537688442211	0.457488371822925\\
0.512562814070352	0.452749154685067\\
0.517587939698492	0.448362110153251\\
0.522613065326633	0.444326357232992\\
0.527638190954774	0.440637426494787\\
0.532663316582915	0.437286817393034\\
0.537688442211055	0.434261640233456\\
0.542713567839196	0.431544371744105\\
0.547738693467337	0.429112745944032\\
0.552763819095477	0.426939791074202\\
0.557788944723618	0.424994010201422\\
0.562814070351759	0.42323969002598\\
0.5678391959799	0.421637312181338\\
0.57286432160804	0.420144036514321\\
0.577889447236181	0.418714228204724\\
0.582914572864322	0.417300010373097\\
0.587939698492462	0.415851839514471\\
0.592964824120603	0.41431911953469\\
0.597989949748744	0.412650887132143\\
0.603015075376884	0.410796612308204\\
0.608040201005025	0.40870715914832\\
0.613065326633166	0.406335941445298\\
0.618090452261307	0.403640285046013\\
0.623115577889447	0.400582976060671\\
0.628140703517588	0.39713393544438\\
0.633165829145729	0.393271921693209\\
0.638190954773869	0.388986131011171\\
0.64321608040201	0.384277544558943\\
0.648241206030151	0.379159870169372\\
0.653266331658292	0.37365994369981\\
0.658291457286432	0.367817492404749\\
0.663316582914573	0.36168421553767\\
0.668341708542714	0.355322199247343\\
0.673366834170854	0.34880174543529\\
0.678391959798995	0.342198749053744\\
0.683417085427136	0.335591798099582\\
0.688442211055276	0.329059190601626\\
0.693467336683417	0.322676061764477\\
0.698492462311558	0.316511793947012\\
0.703517587939699	0.310627846771036\\
0.708542713567839	0.305076100371572\\
0.71356783919598	0.299897757914833\\
0.718592964824121	0.295122809527867\\
0.723618090452261	0.290770022669268\\
0.728643216080402	0.286847395847397\\
0.733668341708543	0.283352993857986\\
0.738693467336683	0.280276072433926\\
0.743718592964824	0.277598396626824\\
0.748743718592965	0.275295658274199\\
0.753768844221106	0.273338901406472\\
0.758793969849246	0.271695868318863\\
0.763819095477387	0.270332181091915\\
0.768844221105528	0.269212270988396\\
0.773869346733668	0.268299957874232\\
0.778894472361809	0.267558558529375\\
0.78391959798995	0.266950358644305\\
0.78894472361809	0.266435210543229\\
0.793969849246231	0.265967934450781\\
0.798994974874372	0.265493977703159\\
0.804020100502513	0.26494240097941\\
0.809045226130653	0.264215421214477\\
0.814070351758794	0.263172781354016\\
0.819095477386935	0.261609528275579\\
0.824120603015075	0.259226513826231\\
0.829145728643216	0.255596259760192\\
0.834170854271357	0.250135754644893\\
0.839195979899497	0.242116882831704\\
0.844221105527638	0.230776530282181\\
0.849246231155779	0.215615771522519\\
0.85427135678392	0.196935161468144\\
0.85929648241206	0.176418538052175\\
0.864321608040201	0.157138822882367\\
0.869346733668342	0.142265183930238\\
0.874371859296482	0.133054340911765\\
0.879396984924623	0.1283883066612\\
0.884422110552764	0.126342380207929\\
0.889447236180904	0.12551645367536\\
0.894472361809045	0.12519610906771\\
0.899497487437186	0.125074034993859\\
0.904522613065327	0.125027870047439\\
0.909547738693467	0.125010473223589\\
0.914572864321608	0.12500393019789\\
0.919597989949749	0.125001472836433\\
0.924623115577889	0.125000551136018\\
0.92964824120603	0.125000205901365\\
0.934673366834171	0.12500007678448\\
0.939698492462312	0.125000028576527\\
0.944723618090452	0.125000010611276\\
0.949748743718593	0.125000003930448\\
0.954773869346734	0.125000001451848\\
0.959798994974874	0.125000000534674\\
0.964824120603015	0.125000000196257\\
0.969849246231156	0.125000000071781\\
0.974874371859296	0.125000000026153\\
0.979899497487437	0.125000000009489\\
0.984924623115578	0.125000000003428\\
0.989949748743719	0.125000000001232\\
0.994974874371859	0.125000000000441\\
1	0.125000000000441\\
};
\addlegendentry{2nd Order Temporal}

\end{axis}
\end{tikzpicture}%}
    };
    
    \node[anchor=south west, fill=white] at (\hspace,\vspace) {
        \resizebox{\plotwidth}{\plotheight}{% This file was created by matlab2tikz.
%
%The latest updates can be retrieved from
%  http://www.mathworks.com/matlabcentral/fileexchange/22022-matlab2tikz-matlab2tikz
%where you can also make suggestions and rate matlab2tikz.
%
\definecolor{mycolor1}{rgb}{0.12941,0.12941,0.12941}%
%
\begin{tikzpicture}

\begin{axis}[%
width=3.229in,
height=2.498in,
at={(0.542in,0.392in)},
scale only axis,
xmin=0,
xmax=1,
xlabel style={font=\color{mycolor1}},
xlabel={x},
ymin=0,
ymax=1,
ylabel style={font=\color{mycolor1}},
ylabel={Velocity (u)},
axis background/.style={fill=white},
title style={font=\bfseries\color{mycolor1}},
title={Second-Order Temporal Velocity Distribution},
xmajorgrids,
ymajorgrids,
legend style={legend cell align=left, align=left}
]
\addplot [color=red, line width=2.0pt]
  table[row sep=crcr]{%
0	4.96997182411239e-12\\
0.0050251256281407	4.96997182411239e-12\\
0.0100502512562814	1.09124957536652e-11\\
0.0150753768844221	2.36851556667719e-11\\
0.0201005025125628	5.08138185605085e-11\\
0.0251256281407035	1.07746918677238e-10\\
0.0301507537688442	2.25793249078473e-10\\
0.0351758793969849	4.67589107054701e-10\\
0.0402010050251256	9.56815914159698e-10\\
0.0452261306532663	1.9344852049519e-09\\
0.050251256281407	3.86399731734005e-09\\
0.0552763819095477	7.62435081880798e-09\\
0.0603015075376884	1.48601364419636e-08\\
0.0653266331658292	2.86058399690955e-08\\
0.0703517587939698	5.4382048400748e-08\\
0.0753768844221105	1.02089517812471e-07\\
0.0804020100502513	1.89228280858979e-07\\
0.085427135678392	3.46276922952599e-07\\
0.0904522613065327	6.25526475939251e-07\\
0.0954773869346734	1.11532949787192e-06\\
0.100502512562814	1.9626613979503e-06\\
0.105527638190955	3.40816221490487e-06\\
0.110552763819095	5.83948335236697e-06\\
0.115577889447236	9.87081760879948e-06\\
0.120603015075377	1.64588805988697e-05\\
0.125628140703518	2.70681540481185e-05\\
0.130653266331658	4.39005360588493e-05\\
0.135678391959799	7.02060796280854e-05\\
0.14070351758794	0.000110691385491721\\
0.14572864321608	0.000172039350990939\\
0.150753768844221	0.000263547123820248\\
0.155778894472362	0.000397877113234509\\
0.160804020100503	0.000591898071083777\\
0.165829145728643	0.000867569823524847\\
0.170854271356784	0.00125279796237877\\
0.175879396984925	0.00178215736703641\\
0.180904522613065	0.00249736147939146\\
0.185929648241206	0.00344734485665639\\
0.190954773869347	0.0046878367619072\\
0.195979899497487	0.00628033840865965\\
0.201005025125628	0.00829047646788733\\
0.206030150753769	0.0107857847993961\\
0.21105527638191	0.0138330527546184\\
0.21608040201005	0.0174954548862273\\
0.221105527638191	0.0218297260407183\\
0.226130653266332	0.0268836547209243\\
0.231155778894472	0.0326941322560066\\
0.236180904522613	0.0392859222325304\\
0.241206030150754	0.0466712191905048\\
0.246231155778894	0.0548499674376902\\
0.251256281407035	0.0638108283262639\\
0.256281407035176	0.0735326298830794\\
0.261306532663317	0.0839861106325629\\
0.266331658291457	0.0951357766605733\\
0.271356783919598	0.106941719337056\\
0.276381909547739	0.119361280750511\\
0.281407035175879	0.132350495799261\\
0.28643216080402	0.145865277491735\\
0.291457286432161	0.159862341697744\\
0.296482412060302	0.174299888261662\\
0.301507537688442	0.189138067689549\\
0.306532663316583	0.204339268151554\\
0.311557788944724	0.219868258186743\\
0.316582914572864	0.235692218022428\\
0.321608040201005	0.25178068825225\\
0.326633165829146	0.268105459794304\\
0.331658291457286	0.284640424251925\\
0.336683417085427	0.301361399418641\\
0.341708542713568	0.318245940892154\\
0.346733668341709	0.335273147641311\\
0.351758793969849	0.352423466875431\\
0.35678391959799	0.369678501626433\\
0.361809045226131	0.3870208229845\\
0.366834170854271	0.40443378784105\\
0.371859296482412	0.421901362210366\\
0.376884422110553	0.439407949657248\\
0.381909547738693	0.456938223999487\\
0.386934673366834	0.474476965240811\\
0.391959798994975	0.492008897593682\\
0.396984924623116	0.509518528454117\\
0.402010050251256	0.526989987284049\\
0.407035175879397	0.544406863540598\\
0.412060301507538	0.561752043073739\\
0.417085427135678	0.579007542809083\\
0.422110552763819	0.596154344062422\\
0.42713567839196	0.613172225524581\\
0.4321608040201	0.630039597840844\\
0.437185929648241	0.646733342822973\\
0.442211055276382	0.663228661706173\\
0.447236180904523	0.679498938524017\\
0.452261306532663	0.695515626630712\\
0.457286432160804	0.711248168633983\\
0.462311557788945	0.726663962450179\\
0.467336683417085	0.741728388728608\\
0.472361809045226	0.756404917298663\\
0.477386934673367	0.770655312244071\\
0.482412060301508	0.784439956247001\\
0.487437185929648	0.797718314381461\\
0.492462311557789	0.810449554870869\\
0.49748743718593	0.822593338714039\\
0.50251256281407	0.83411078085623\\
0.507537688442211	0.844965572324093\\
0.512562814070352	0.855125235537115\\
0.517587939698492	0.864562464690398\\
0.522613065326633	0.873256481481142\\
0.527638190954774	0.881194316374273\\
0.532663316582915	0.888371910788569\\
0.537688442211055	0.894794930109131\\
0.542713567839196	0.900479184889513\\
0.547738693467337	0.905450580065503\\
0.552763819095477	0.909744549024509\\
0.557788944723618	0.913404977453916\\
0.562814070351759	0.91648267463117\\
0.5678391959799	0.919033498972782\\
0.57286432160804	0.921116281879938\\
0.577889447236181	0.922790712728142\\
0.582914572864322	0.924115345242465\\
0.587939698492462	0.925145862644297\\
0.592964824120603	0.925933700768808\\
0.597989949748744	0.926525082146439\\
0.603015075376884	0.926960467695184\\
0.608040201005025	0.927274392912292\\
0.613065326633166	0.927495626707993\\
0.618090452261307	0.927647575009192\\
0.623115577889447	0.92774884728007\\
0.628140703517588	0.927813909772679\\
0.633165829145729	0.927853761433013\\
0.638190954773869	0.927876583707806\\
0.64321608040201	0.927888331289154\\
0.648241206030151	0.927893245138154\\
0.653266331658292	0.927894280804057\\
0.658291457286432	0.927893453669738\\
0.663316582914573	0.927892108398757\\
0.668341708542714	0.927891122950262\\
0.673366834170854	0.927891058633678\\
0.678391959798995	0.927892267386083\\
0.683417085427136	0.927894966300957\\
0.688442211055276	0.927899287842318\\
0.693467336683417	0.927905312448622\\
0.698492462311558	0.927913088558147\\
0.703517587939699	0.92792264356372\\
0.708542713567839	0.92793398783911\\
0.71356783919598	0.927947112711348\\
0.718592964824121	0.927961981956492\\
0.723618090452261	0.927978514870702\\
0.728643216080402	0.927996556914132\\
0.733668341708543	0.928015830892559\\
0.738693467336683	0.928035856948906\\
0.743718592964824	0.928055822237064\\
0.748743718592965	0.928074369428279\\
0.753768844221106	0.928089254657204\\
0.758793969849246	0.928096796283756\\
0.763819095477387	0.928090989942079\\
0.768844221105528	0.928062093484905\\
0.773869346733668	0.9279943732288\\
0.778894472361809	0.927862528076352\\
0.78391959798995	0.927626036102576\\
0.78894472361809	0.92722024716661\\
0.793969849246231	0.926542526159413\\
0.798994974874372	0.925430798151034\\
0.804020100502513	0.923629867558791\\
0.809045226130653	0.920740009409285\\
0.814070351758794	0.916138933853159\\
0.819095477386935	0.908865914910348\\
0.824120603015075	0.897455773862636\\
0.829145728643216	0.879715866610526\\
0.834170854271357	0.852463208458056\\
0.839195979899497	0.811307317040352\\
0.844221105527638	0.750724041061721\\
0.849246231155779	0.664971695041896\\
0.85427135678392	0.550800346327225\\
0.85929648241206	0.412759007668186\\
0.864321608040201	0.268990499347793\\
0.869346733668342	0.148051767307356\\
0.874371859296482	0.0693839450513218\\
0.879396984924623	0.0290365580172406\\
0.884422110552764	0.0114382115146456\\
0.889447236180904	0.00438537238882011\\
0.894472361809045	0.0016624101521941\\
0.899497487437186	0.000627127617386477\\
0.904522613065327	0.000236006338842851\\
0.909547738693467	8.86774517682004e-05\\
0.914572864321608	3.3275641679873e-05\\
0.919597989949749	1.24697714677805e-05\\
0.924623115577889	4.66616063841524e-06\\
0.92964824120603	1.74324705714763e-06\\
0.934673366834171	6.50088846657838e-07\\
0.939698492462312	2.41940482357912e-07\\
0.944723618090452	8.98393622437208e-08\\
0.949748743718593	3.3276758512151e-08\\
0.954773869346734	1.22919308606243e-08\\
0.959798994974874	4.52676409907143e-09\\
0.964824120603015	1.66159259548183e-09\\
0.969849246231156	6.07725911605549e-10\\
0.974874371859296	2.21418336868944e-10\\
0.979899497487437	8.03372481620404e-11\\
0.984924623115578	2.90196917667106e-11\\
0.989949748743719	1.04331896241398e-11\\
0.994974874371859	3.7321943442894e-12\\
1	3.7321943442894e-12\\
};
\addlegendentry{2nd Order Temporal}

\end{axis}
\end{tikzpicture}%}
    };
    
    % Bottom row
    \node[anchor=south west, fill=white] at (0,0) {
        \resizebox{\plotwidth}{\plotheight}{% This file was created by matlab2tikz.
%
%The latest updates can be retrieved from
%  http://www.mathworks.com/matlabcentral/fileexchange/22022-matlab2tikz-matlab2tikz
%where you can also make suggestions and rate matlab2tikz.
%
\definecolor{mycolor1}{rgb}{0.12941,0.12941,0.12941}%
%
\begin{tikzpicture}

\begin{axis}[%
width=3.229in,
height=2.498in,
at={(0.542in,0.392in)},
scale only axis,
xmin=0,
xmax=1,
xlabel style={font=\color{mycolor1}},
xlabel={x},
ymin=0.1,
ymax=1,
ylabel style={font=\color{mycolor1}},
ylabel={Pressure (p)},
axis background/.style={fill=white},
title style={font=\bfseries\color{mycolor1}},
title={Second-Order Temporal Pressure Distribution},
xmajorgrids,
ymajorgrids,
legend style={legend cell align=left, align=left}
]
\addplot [color=green, line width=2.0pt]
  table[row sep=crcr]{%
0	0.999999999994119\\
0.0050251256281407	0.999999999994119\\
0.0100502512562814	0.999999999987088\\
0.0150753768844221	0.999999999971975\\
0.0201005025125628	0.999999999939876\\
0.0251256281407035	0.999999999872512\\
0.0301507537688442	0.999999999732838\\
0.0351758793969849	0.999999999446741\\
0.0402010050251256	0.99999999886788\\
0.0452261306532663	0.999999997711086\\
0.050251256281407	0.999999995428057\\
0.0552763819095477	0.999999990978747\\
0.0603015075376884	0.99999998241725\\
0.0653266331658292	0.999999966153114\\
0.0703517587939698	0.999999935654294\\
0.0753768844221105	0.999999879206059\\
0.0804020100502513	0.999999776102098\\
0.085427135678392	0.999999590279684\\
0.0904522613065327	0.999999259867305\\
0.0954773869346734	0.999998680325019\\
0.100502512562814	0.999997677749822\\
0.105527638190955	0.999995967414455\\
0.110552763819095	0.999993090648885\\
0.115577889447236	0.999988320744904\\
0.120603015075377	0.999980525739939\\
0.125628140703518	0.999967972935503\\
0.130653266331658	0.999948057260041\\
0.135678391959799	0.999916933803885\\
0.14070351758794	0.999869035062797\\
0.14572864321608	0.99979645695639\\
0.150753768844221	0.999688206055412\\
0.155778894472362	0.999529315097596\\
0.160804020100503	0.99929985571883\\
0.165829145728643	0.998973906172509\\
0.170854271356784	0.998518565639393\\
0.175879396984925	0.997893141134454\\
0.180904522613065	0.997048661024605\\
0.185929648241206	0.995927881836608\\
0.190954773869347	0.994465943057562\\
0.195979899497487	0.992591780995966\\
0.201005025125628	0.990230335919732\\
0.206030150753769	0.987305483192998\\
0.21105527638191	0.983743504547888\\
0.21608040201005	0.979476812389281\\
0.221105527638191	0.974447572279353\\
0.226130653266332	0.968610855288512\\
0.231155778894472	0.961936999771964\\
0.236180904522613	0.954412963486503\\
0.241206030150754	0.946042580812792\\
0.246231155778894	0.936845778240939\\
0.251256281407035	0.926856917487795\\
0.256281407035176	0.916122511210323\\
0.261306532663317	0.904698584657021\\
0.266331658291457	0.892647942363958\\
0.271356783919598	0.880037554276628\\
0.276381909547739	0.866936215303776\\
0.281407035175879	0.853412569733596\\
0.28643216080402	0.839533536799339\\
0.291457286432161	0.825363131136472\\
0.296482412060302	0.810961643251341\\
0.301507537688442	0.796385129128918\\
0.306532663316583	0.781685152104182\\
0.311557788944724	0.766908721116029\\
0.316582914572864	0.752098374736385\\
0.321608040201005	0.737292367770942\\
0.326633165829146	0.722524925242594\\
0.331658291457286	0.707826536235833\\
0.336683417085427	0.69322426687799\\
0.341708542713568	0.678742077445553\\
0.346733668341709	0.664401133189603\\
0.351758793969849	0.650220102062692\\
0.35678391959799	0.636215435240351\\
0.361809045226131	0.622401628318457\\
0.366834170854271	0.608791462480411\\
0.371859296482412	0.595396225895293\\
0.376884422110553	0.5822259162382\\
0.381909547738693	0.56928942560376\\
0.386934673366834	0.556594709280576\\
0.391959798994975	0.544148939918848\\
0.396984924623116	0.531958648592329\\
0.402010050251256	0.52002985415479\\
0.407035175879397	0.508368182137083\\
0.412060301507538	0.496978974233443\\
0.417085427135678	0.48586738918925\\
0.422110552763819	0.475038495627176\\
0.42713567839196	0.464497357031645\\
0.4321608040201	0.454249108747826\\
0.437185929648241	0.444299026435302\\
0.442211055276382	0.434652584943111\\
0.447236180904523	0.425315506039595\\
0.452261306532663	0.41629379284028\\
0.457286432160804	0.407593748141365\\
0.462311557788945	0.399221973209722\\
0.467336683417085	0.39118534294508\\
0.472361809045226	0.383490952782684\\
0.477386934673367	0.376146032340219\\
0.482412060301508	0.369157820759212\\
0.487437185929648	0.362533399108167\\
0.492462311557789	0.356279476287102\\
0.49748743718593	0.350402126791749\\
0.50251256281407	0.34490648162564\\
0.507537688442211	0.339796377680428\\
0.512562814070352	0.335073975992758\\
0.517587939698492	0.330739365177478\\
0.522613065326633	0.326790172513942\\
0.527638190954774	0.323221210816552\\
0.532663316582915	0.320024193292899\\
0.537688442211055	0.317187549900098\\
0.542713567839196	0.314696376169746\\
0.547738693467337	0.312532538397978\\
0.552763819095477	0.310674947499706\\
0.557788944723618	0.309099998623167\\
0.562814070351759	0.307782156653314\\
0.5678391959799	0.306694651523546\\
0.57286432160804	0.305810234524317\\
0.577889447236181	0.305101939841326\\
0.582914572864322	0.304543795675717\\
0.587939698492462	0.304111436455486\\
0.592964824120603	0.303782580471057\\
0.597989949748744	0.303537353411798\\
0.603015075376884	0.30335845500279\\
0.608040201005025	0.303231180729467\\
0.613065326633166	0.303143321689664\\
0.618090452261307	0.303084972072086\\
0.623115577889447	0.303048275690314\\
0.628140703517588	0.303027141163172\\
0.633165829145729	0.303016950871639\\
0.638190954773869	0.303014282966836\\
0.64321608040201	0.303016659532494\\
0.648241206030151	0.303022328321586\\
0.653266331658292	0.303030080781489\\
0.658291457286432	0.303039105564303\\
0.663316582914573	0.303048874382559\\
0.668341708542714	0.303059055769916\\
0.673366834170854	0.303069451824691\\
0.678391959798995	0.303079953116275\\
0.683417085427136	0.303090507402707\\
0.688442211055276	0.303101098460637\\
0.693467336683417	0.303111732029264\\
0.698492462311558	0.303122426522587\\
0.703517587939699	0.30313320671012\\
0.708542713567839	0.303144098970301\\
0.71356783919598	0.303155126961883\\
0.718592964824121	0.303166306614525\\
0.723618090452261	0.303177639174529\\
0.728643216080402	0.303189100587277\\
0.733668341708543	0.303200624633847\\
0.738693467336683	0.303212075761249\\
0.743718592964824	0.303223205121376\\
0.748743718592965	0.303233579434505\\
0.753768844221106	0.303242466090329\\
0.758793969849246	0.303248648104759\\
0.763819095477387	0.303250127170113\\
0.768844221105528	0.30324364899053\\
0.773869346733668	0.303223947652152\\
0.778894472361809	0.303182547711027\\
0.78391959798995	0.303105873052629\\
0.78894472361809	0.302972274584187\\
0.793969849246231	0.302747425894915\\
0.798994974874372	0.302377243752335\\
0.804020100502513	0.301776896729987\\
0.809045226130653	0.300814339773819\\
0.814070351758794	0.299286082459927\\
0.819095477386935	0.296883066371101\\
0.824120603015075	0.293146254681608\\
0.829145728643216	0.287417662714676\\
0.834170854271357	0.278808139077453\\
0.839195979899497	0.266234956226278\\
0.844221105527638	0.248630065282384\\
0.849246231155779	0.225448788691989\\
0.85427135678392	0.197502425961491\\
0.85929648241206	0.167721432470044\\
0.864321608040201	0.140828667672924\\
0.869346733668342	0.12107765860296\\
0.874371859296482	0.109487728603911\\
0.879396984924623	0.103892911493847\\
0.884422110552764	0.101520944377586\\
0.889447236180904	0.100581262727308\\
0.894472361809045	0.100220077906106\\
0.899497487437186	0.100082984132523\\
0.904522613065327	0.100031223948997\\
0.909547738693467	0.100011731381911\\
0.914572864321608	0.100004402018027\\
0.919597989949749	0.100001649604747\\
0.924623115577889	0.100000617276296\\
0.92964824120603	0.100000230610086\\
0.934673366834171	0.100000085998695\\
0.939698492462312	0.100000032005721\\
0.944723618090452	0.100000011884631\\
0.949748743718593	0.100000004402101\\
0.954773869346734	0.10000000162607\\
0.959798994974874	0.100000000598835\\
0.964824120603015	0.100000000219808\\
0.969849246231156	0.100000000080395\\
0.974874371859296	0.100000000029291\\
0.979899497487437	0.100000000010628\\
0.984924623115578	0.100000000003839\\
0.989949748743719	0.10000000000138\\
0.994974874371859	0.100000000000494\\
1	0.100000000000494\\
};
\addlegendentry{2nd Order Temporal}

\end{axis}
\end{tikzpicture}%}
    };
    
\end{tikzpicture}
\caption{Kết quả phương pháp bậc 2 thời gian tại t = 0.2, kết hợp công thức minmod}
\label{fig:temporal2_results}
\end{figure}
